\documentclass[11pt,a4paper]{article}

% --- Packages ---
\usepackage[utf8]{inputenc}
\usepackage[T1]{fontenc}
\usepackage{amsmath,amssymb,amsthm}
\usepackage{booktabs}
\usepackage{graphicx}
\usepackage{xurl}
\usepackage{hyperref}
\usepackage{cleveref}
\usepackage{xcolor}
\usepackage{enumitem}
\usepackage{microtype}
\usepackage{natbib}
\usepackage{authblk}
\usepackage[margin=1in]{geometry}
\usepackage{algorithm}
\usepackage{algpseudocode}
\usepackage{multirow}
\usepackage{subcaption}
\usepackage{tabularx}

% --- Theorem environments ---
\newtheorem{theorem}{Theorem}
\newtheorem{lemma}[theorem]{Lemma}
\newtheorem{proposition}[theorem]{Proposition}
\newtheorem{corollary}[theorem]{Corollary}
\newtheorem{definition}{Definition}

% --- Macros ---
\newcommand{\R}{\mathbb{R}}
\newcommand{\E}{\mathbb{E}}
\newcommand{\N}{\mathbb{N}}
\newcommand{\bx}{\mathbf{x}}
\newcommand{\bh}{\mathbf{h}}
\newcommand{\bW}{\mathbf{W}}
\newcommand{\bQ}{\mathbf{Q}}
\newcommand{\bK}{\mathbf{K}}
\newcommand{\bV}{\mathbf{V}}
\newcommand{\bg}{\mathbf{g}}
\newcommand{\bmask}{\mathbf{m}}
\newcommand{\bz}{\mathbf{z}}
\newcommand{\cR}{\mathcal{R}}
\newcommand{\cL}{\mathcal{L}}

\hypersetup{
    pdftitle={Cryptographic Dimension Partitioning: FFN Confinement and Whole-Model Authorization},
    pdfauthor={Ryan McCormick},
    pdfsubject={Identity-bound FFN confinement and whole-model authorization},
    pdfkeywords={cryptographic dimension partitioning, FFN gating, whole-model authorization, multi-tenant neural networks, Lean 4},
    colorlinks=true,
    linkcolor=blue!70!black,
    citecolor=blue!70!black,
    urlcolor=blue!70!black
}

% ===================================================================
\title{\textbf{Cryptographic Dimension Partitioning} \\
\large{Identity-Bound FFN Confinement and Whole-Model Authorization}}

\author[1]{Ryan McCormick}
\affil[1]{Independent Researcher}
\date{August 2026}

\begin{document}
\maketitle

% ===================================================================
\begin{abstract}
Public-key infrastructure made identity administrable: certificates
bind keys to names, and signed dereferences compose across
organizations that share no other trust.  Model artifacts need an
analogous infrastructure for \emph{capability}.  As trained models are
licensed, co-located, and embedded in other parties' systems, what is
required is not merely encryption of the artifact but an
identity-bound structure \emph{inside} it to which signed grants of
scoped capability can unambiguously refer --- a public-key
infrastructure for model capability.  This paper develops the
structural basis for that infrastructure.
Shared multi-tenant systems lack a structural zero baseline for
tenant-scoped learned state and often leave complete-model dispatch to
serving configuration.  We introduce \emph{Cryptographic Dimension
Partitioning} (CDP), which binds an authenticated tenant + sub-id scope
to two distinct gate surfaces.  At an FFN boundary, a runtime-selected
binary mask may act on the pre-MLP input or on the expanded
intermediate activation.  Both placements give exact zeros at the
gated tensor; the intermediate placement additionally aligns the
hidden columns of $\bW_1$, $\mathbf b_1$, and rows of $\bW_2$ for the
paper's strongest parameter-confinement result.  At the whole-model
surface, a pre-execution authorization gate treats the complete model
as one atomic resource: an authorized request executes the unchanged
model, while a rejected request never invokes it.

We verify the gate, gradient, and aligned FFN confinement results in
Lean~4 with no \texttt{sorry} declarations or project-specific axioms
in the headline theorems.  A complementary negative theorem shows
that placing LayerNorm between the activation and gate breaks
$\bW_1$ confinement, making placement and path closure explicit
preconditions.

The guarantee classes are of different strength, and the empirics are
stated in matching strength.  Structurally, inactive-state exclusion is exact:
in a strict frozen-backbone DistilBERT experiment, every intermediate
FFN is gated and only authorized FFN slices plus a private classifier
are trainable, and across three seeds the maximum frozen-shared,
off-partition parameter, off-partition optimizer-moment, and
inactive-classifier delta is exactly zero in every run.  Where the
network has capacity, partitioning also costs little to nothing
measurably: in a five-seed paired sweep at a post-encoder gate,
co-trained DistilBERT stays within $0.14$--$0.38$ percentage points of
separately trained models across $R{=}8$--$128$ --- no paired test
survives Bonferroni correction --- and in a synthetic capacity-scaling
check at $R{=}128$ with per-regime width held fixed, all regimes train
to $100\%$ owning accuracy, zero of $32{,}512$ wrong-key probes are
answered, and support extraction is exact (maximum logit difference
$0$).  Where capacity is genuinely scarce the cost is real and
recoverable: strict frozen-backbone owning accuracy decreases from
$91.28\%\pm0.33\%$ at $R{=}1$ to $86.21\%\pm0.08\%$ at $R{=}16$,
and complete private residual adapters and private experts recover
$90.47\%$ and $90.09\%$ mean accuracy,
respectively, without modifying the shared backbone or inactive
lanes. In an $R{=}8$ learned-MoE study on 6,231 held-out examples,
separately trained top-1 routers and experts retain exactly equal logits and
predictions after packing into one authorized bank, with zero unauthorized
dispatch. A token-level extension makes 311,477 held-out routing decisions
under execution-selected private continuous prefixes; literal capability-marker
text cannot expand authority, and unauthorized dispatch remains zero. The
source-visible research preflight checks pre-MLP exclusion and fail-closed
whole-model invocation separately.  The completed results support the
declared FFN-coordinate, private-state, and atomic model-authorization
boundaries; they do not establish multi-regime partitioning inside
shared attention.\footnote{Code, artifacts,
and proofs: \url{https://github.com/sekosai/schemen-gate/tree/main/research/cdp}}
\end{abstract}

\vspace{0.5em}
\noindent\textbf{Keywords:} multi-tenant neural networks,
cryptographic isolation, dimension partitioning, gate masks,
Hadamard gating, formal verification, Lean~4, support-constrained
training, permutation equivariance, classification, submodel extraction

% ===================================================================
\section{Introduction}
\label{sec:intro}

\paragraph{An identity infrastructure for model capability.}
X.509-style public-key infrastructure answered a precise question:
\emph{whose key is this, and who vouches for it?}  A certificate
binding a key to a name lets signed authority traverse organizations
without shared custody.  As trained models move from vendor servers
into customer laptops, partner runtimes, and regulated estates, the
analogous question about model artifacts becomes operational: \emph{of
the capability inside this artifact, which scope is this party
entitled to execute, and whose signature says so?}  Answering it
requires more than authenticating a channel or encrypting a file: it
requires identity-bound structure inside the artifact that signed
grants can address the way certificates address keys --- the internal
capability surface of a public-key infrastructure for model capability.
Cryptographic Dimension Partitioning (CDP) is developed here as that surface.

The economics of large language model (LLM) deployment push
strongly toward sharing.  A single 7-billion-parameter model
requires roughly 14\,GB for float16 weights or 7\,GB for int8 weights,
before KV caches, activations, allocator overhead, and serving state.
Serving $N$ tenants on $N$ separate model copies
scales linearly in cost.  Parameter-efficient fine-tuning methods
such as LoRA~\citep{hu2022lora} offer an alternative: train a
small adapter per tenant while sharing the backbone.  Platforms
like Predibase~\citep{predibase2024lorax} and
S-LoRA~\citep{sheng2024slora} serve hundreds of adapters from a
single GPU, promising multi-tenant efficiency with per-tenant
customization.

LoRA adds a low-rank perturbation $\mathbf{B}\mathbf{A}$ to selected
weight matrices while the base weights are commonly frozen.  Its
multi-tenant serving boundary is the adapter router and the associated
batch, cache, and lifecycle state.  A correctly routed adapter changes
the hidden representation of its own request; that request-local
diffusion is not, by itself, cross-tenant leakage.  Another tenant is
exposed only if routing, batching, cache namespacing, shared training,
or another serving boundary admits the wrong adapter state.  LoRA
therefore provides efficient private state but does not inherently
supply an IdP-bound cryptographic authorization gate or a
coordinate-zero invariant.

Common LoRA configurations target attention projections ($\bW_Q$,
$\bW_K$, $\bW_V$, $\bW_O$), FFN layers, or both.  Inside one active
request, a dense shared attention head cannot be subdivided into
independently normalized coordinate channels merely by assigning
subsets inside its dot product: the partial sums are combined into one
score before the shared softmax.  This is a limitation on
intra-attention coordinate partitioning, not a claim that a properly
routed LoRA adapter is visible to every tenant.

\begin{proposition}[Dense within-head attention coupling]
\label{prop:attention-entanglement}
Let $\bW_Q' = \bW_Q + \mathbf{B}_Q\mathbf{A}_Q$ be a LoRA
modification to the query projection, where
$\mathbf{B}_Q \in \R^{d \times r}$,
$\mathbf{A}_Q \in \R^{r \times d_k}$.  For any partition
$\{1, \ldots, d_k\} = S_i \sqcup S_j$, the modified attention
score algebraically decomposes as
\[
    (\bQ'\bK^\top)_{p,q} = \sum_{m=1}^{d_k} Q'_{p,m} K_{q,m}
    = \underbrace{\sum_{m \in S_i} Q'_{p,m} K_{q,m}}_{\text{subset } i}
    + \underbrace{\sum_{m \in S_j} Q'_{p,m} K_{q,m}}_{\text{subset } j}
\]
but this arithmetic decomposition does not produce independently
owned channels.  A generic dense update
$\mathbf{B}_Q\mathbf{A}_Q$ is not support-constrained and may alter
$Q'_{p,m}$ on either subset.  Even if an update is constrained to one
subset, the two partial scores are added before a single shared
softmax; after that addition, a post-hoc coordinate mask cannot
recover separately normalized tenant scores.  Thus coordinate
partitioning alone does not isolate tenants within one dense shared
head.  The proposition does not rule out a different architecture
that computes complete attention and residual lanes separately and
gates whole lanes before any recombination.
\end{proposition}

CDP has two architecture scopes.  For partitioning within one model,
the gate acts at an FFN boundary, where element-wise masking permits
exact coordinate exclusion at the declared tensor.  A
\emph{pre-MLP} gate masks the FFN input before $\bW_1$; it is still an
FFN placement, but its aligned parameter surface is limited to the
input-facing rows of $\bW_1$.  An \emph{intermediate-FFN} gate masks
the expanded activation before $\bW_2$ and confines the corresponding
$\bW_1$, $\mathbf b_1$, and $\bW_2$ slices under a
coordinate-respecting optimizer.  The latter is the strongest
formally and empirically evaluated placement in this paper.  An
extended attention embodiment is possible, but it is not a
coordinate mask inside an ordinary shared head: it duplicates each
Transformer block into regime-specific lanes, including each lane's
complete set of heads and every residual path, and gates the whole
lane before cross-lane recombination.  That construction escapes the
within-head premise of the proposition, but its duplicated topology
creates a materially larger operator-error surface.  Unless every
attention and residual path is lane-aligned, attention remains inside
the shared trust boundary.  For the core FFN gate,
$\bh = \sigma(\bx \bW_1 + \mathbf{b}_1)$; each hidden dimension
$h_j$ is computed independently.  Zeroing $h_j$ via the gate mask
eliminates dimension $j$'s contribution to the output exactly.
No other hidden coordinate is changed by the multiplication itself;
the subsequent $\bW_2$ may mix the remaining active contributions
across output coordinates.

For whole-model authorization, the complete model is instead one
atomic resource behind a fail-closed runtime decision.  The authorized
path does not ablate or repartition internal activations.  These FFN
and whole-model scopes can be nested, but their guarantees are
different.

Within either FFN placement, the fundamental mechanism is
multiplication by a binary mask:
$\bg_r \in \{0,1\}^d$ is a vector of zeros and ones, and the
gated activation is $\tilde{\bh} = \bh \odot \bg_r$.  Ones pass
signal through; zeros annihilate it.  This is the local confinement
mechanism; authenticated dispatch and closure of other paths define
the larger deployment boundary.

The \emph{security layer} determines how masks are assigned and who
may activate them.  A lockbox-held assignment key drives an
HMAC-SHA256 stream and unbiased Fisher--Yates shuffle, producing one
global partition of the hidden dimension into $R$ disjoint masks.
The lockbox then issues scoped authorization material for an assigned
mask; independently shuffling under each child key would not preserve
global disjointness and is not the protocol.  The authenticated
runtime resolves a principal and operation to the assigned mask, so
callers cannot choose a projection directly.

The resolved regime has a dual security role.  Its selection is an
authorization decision at the ingress boundary, while the resulting
cryptographically bound regime is the authenticated execution identity
for each downstream gate that verifies the same grant, model, operation,
and policy scope.  Thus the chain is external authentication, boundary
authorization, then downstream regime authentication.  A bare regime
integer or activation tensor does not authenticate itself.

\paragraph{Contributions.}
\begin{enumerate}[nosep]
    \item A capability-PKI formulation for model artifacts:
    identity-bound internal structure (partitions) that signed grants
    can address, with lockbox-generated assignment, so scoped
    capability can be conferred, composed, and revoked across parties
    without shared custody of the weights.
    \item The CDP mechanism: a lockbox-generated global assignment
    from HMAC-SHA256 and unbiased Fisher--Yates to disjoint FFN masks,
    plus tenant + sub-id scoped authorization for FFN activation or
    atomic whole-model invocation.
    \item Hadamard governance: one algebraic rule
    ($\tilde{\bh} = \bh \odot \bg_r$) from which forward coordinate
    exclusion, active preservation, gradient isolation, and aligned
    FFN weight confinement follow.
    \item Formal verification in Lean~4: the checked-in proof
    inventory contains 21 modules and 312 active theorem/lemma declarations
    with no active \texttt{sorry} or \texttt{admit} declarations.  Five
    project-specific \texttt{axiom} declarations remain: two opaque
    predicates, one conditional PRF enumeration bound, and two historical
    statistical consequences.  The core Hadamard and aligned FFN confinement
    results do not depend on them.  The two statistical declarations
    (\texttt{camouflage\_indistinguishable},
    \texttt{gradient\_probing\_hard}) are excluded from the paper's claim set
    because their conclusions are not provable in the white-box setting; the
    declarations remain inspectable as historical conditional modules.  The
    deployments they served are secured by the custody layer instead (see
    \S\ref{sec:formal}).
    \item A placement theorem and counterexample identifying the safe
    \emph{activation $\to$ gate $\to\bW_2$} ordering and showing why
    pre-gate LayerNorm invalidates $\bW_1$ confinement.  The narrower
    pre-MLP placement is specified separately so its input-coordinate
    and first-projection guarantee is not confused with the
    intermediate theorem.
    \item A strict three-seed DistilBERT cotenancy study that freezes
    all non-governed shared paths and measures intermediate-FFN,
    complete private-adapter, and private-expert designs.  All
    unauthorized parameter and optimizer-state deltas are exactly
    zero; the dense FFN capacity curve is measured through $R{=}16$.
    \item A Cargo Mode architecture and self-contained strengthened
    protocol runner binding tenant, sub-identity, regime, model digest,
    and operation before customer context can reach a shared frozen
    Transformer, plus a separate whole-model gate before generator
    invocation.
    \item Supplementary support-constrained classification and keyed
    permutation-conjugation studies, explicitly separated from the
    intermediate-FFN isolation claim.
\end{enumerate}

% ===================================================================
\section{Related Work}
\label{sec:related}

\paragraph{Parameter-efficient fine-tuning.}
LoRA~\citep{hu2022lora} introduces low-rank adapters on attention
projections.  QLoRA~\citep{dettmers2024qlora} extends this with
4-bit base models.  Multi-tenant serving systems
(S-LoRA~\citep{sheng2024slora}, Punica~\citep{chen2023punica},
LoRAX~\citep{predibase2024lorax}) batch multiple adapters
on shared GPUs.  These systems optimize adapter serving rather than
providing a cryptographically authorized FFN-coordinate boundary;
LoRA adapters are commonly installed on shared attention
projections.

\paragraph{Privacy in fine-tuning.}
Differentially Private Stochastic Gradient Descent (DP-SGD)~\citep{abadi2016deep} provides $(\varepsilon,\delta)$-differential
privacy bounds on the influence of training examples.  It does not,
by itself, route requests by tenant identity or allocate disjoint
tenant-scoped model state in a serving process.
\citet{wang2023dpfederated} study federated LoRA with DP
guarantees.  These approaches provide statistical privacy
bounds; CDP provides exact local coordinate confinement at its
declared gate (a different guarantee class).

\paragraph{Model partitioning.}
Expert parallelism in Mixture-of-Experts
(MoE)~\citep{fedus2022switch,lepikhin2021gshard} routes tokens
to different FFN experts.  The routing is learned and data-dependent
(including sparse top-$k$ routing); CDP's activation is instead
resolved from an authenticated capability and a fixed assignment.
The goals differ: MoE targets capacity scaling; CDP targets
authorized FFN-coordinate confinement.

\paragraph{Task-specific masks and subnetworks.}
PackNet~\citep{mallya2018packnet} allocates parameters to sequential
tasks through iterative pruning; Piggyback~\citep{mallya2018piggyback}
learns binary weight masks over a fixed network; and Supermasks in
Superposition~\citep{wortsman2020supermasks} retrieves task-specific
subnetworks from a shared fixed base.  These works establish that
masks and disjoint or reusable subnetworks can support multiple tasks.
CDP does not claim binary masking itself as novel.  Its narrower
contribution is the combination of a globally disjoint FFN activation
assignment, exact aligned optimizer-state confinement, formal
placement conditions, and identity-bound runtime activation.

\paragraph{Confidential and verifiable inference.}
Slalom~\citep{tramer2018slalom} partitions neural-network execution
between trusted hardware and an untrusted accelerator to provide
privacy and integrity for outsourced inference.  CDP assumes rather
than replaces a trusted custody and runtime boundary; its theorem
addresses coordinate support inside that boundary, not confidentiality
against a compromised host or accelerator.

\paragraph{Formal verification of ML systems.}
\citet{seshia2018formal} survey formal methods for ML.  Most
work targets input-output robustness (e.g., certified
adversarial bounds~\citep{cohen2019certified}).  CDP's Lean~4
proofs verify an \emph{architectural} property (gradient and
weight confinement) rather than an input-output property.

\paragraph{Representation superposition.}
\citet{elhage2022superposition} demonstrate that neural networks
encode more features than they have dimensions via nearly-orthogonal
directions.  CDP makes no claim about that microscopic geometry; its
support constraint is simply that excluded coordinates contribute
zero, so each regime's usable gated contribution must be represented
on its active support.

\paragraph{Permutation equivariance.}
The equivariance of linear algebra under change of basis is
elementary.  Our supplementary study applies a keyed residual-stream
permutation consistently to a DistilBERT implementation and validates
function preservation up to $R{=}1{,}024$.  We do not claim the
underlying permutation reparameterization as a new algebraic result.

% ===================================================================
\section{Cryptographic Dimension Partitioning}
\label{sec:cdp}

\subsection{Gate Mask Derivation}
\label{sec:gate-derivation}

\begin{definition}[Gate mask partition]
\label{def:gate-partition}
Given an assignment key $k_{\mathrm{assign}} \in \{0,1\}^{256}$,
a number of regimes
$R \in \N$, and a hidden dimension $d \in \N$ with $d \geq R$,
a \emph{gate mask partition} is a set of binary vectors
$\{\bg_0, \ldots, \bg_{R-1}\} \subset \{0,1\}^d$ satisfying:
\begin{equation}
\bg_r \odot \bg_s = \mathbf{0} \quad \forall\, r \neq s, \qquad
\sum_{r=0}^{R-1} \bg_r = \mathbf{1}
\label{eq:mask-properties}
\end{equation}
where $\odot$ denotes the Hadamard (element-wise) product.  The
masks are disjoint and exhaustive: every dimension is owned by
exactly one regime.  Equal-sized regime partitions additionally
require $R \mid d$, equivalently $d = qR$ for some integer
$q \geq 1$; otherwise, regime sizes differ by at most one.  The
equal-sized construction and capacity formulas used below assume
$R \mid d$ and write $d_r := d/R$ for the dimensions owned by one
regime.
\end{definition}

\begin{algorithm}[ht]
\caption{Global gate-partition assignment with unbiased draws}
\label{alg:gate-mask}
\begin{algorithmic}[1]
\Require Lockbox assignment key $k_{\mathrm{assign}} \in
\{0,1\}^{256}$, model digest $m$, version $v$, number of regimes
$R$, hidden dimension $d$, with $d \geq R$ and $R \mid d$
\Ensure $R$ disjoint binary masks
$\{\bg_0, \ldots, \bg_{R-1}\}$ satisfying~\eqref{eq:mask-properties}
\State Initialize index array $\pi = [0, 1, \ldots, d-1]$
\For{$i = d-1$ down to $1$} \Comment{Fisher-Yates shuffle}
    \State $a \gets 0$; $b \gets i+1$;
    $L \gets 2^{256}-(2^{256}\bmod b)$
    \Repeat
        \State $u \gets \operatorname{OS2IP}(\operatorname{HMAC\mbox{-}SHA256}(
        k_{\mathrm{assign}},
        \operatorname{Encode}(\texttt{partition-v1},m,v,R,d,i,a)))$
        \State $a \gets a+1$
    \Until{$u < L$} \Comment{Reject the incomplete residue tail}
    \State $j \gets u \bmod b$
    \State Swap $\pi[i] \leftrightarrow \pi[j]$
\EndFor
\State Partition $\pi$ into $R$ contiguous blocks of size
$d_r = d/R$
\For{$r = 0$ to $R-1$}
    \State $\bg_r \in \{0,1\}^d$ with $g_{r,j} = 1$ iff
    $j \in \text{block}_r(\pi)$
\EndFor
\end{algorithmic}
\end{algorithm}

The rejection threshold removes the incomplete residue tail before
the modulo operation; \texttt{rejection\_sampling\_count} proves the
corresponding equal-cardinality fact in Lean~4.  Under the standard
HMAC pseudorandom-function (PRF) assumption, the accepted words are
computationally indistinguishable from uniform words.  Rejection
sampling prevents the modulo operation from introducing an additional
residue imbalance into that pseudorandom stream.  The result is deterministic for
the same assignment key and encoded context.  Mask secrecy is defense
in depth: authorization ultimately depends on authenticated dispatch,
not on treating a guessed coordinate set as a credential.

\subsection{Key Hierarchy and Token System}
\label{sec:key-hierarchy}

CDP domain-separates global assignment from authorization.  The
lockbox derives one assignment key for the complete model partition
and separate scoped capability keys:

\begin{align}
k_{\mathrm{assign}} &= \operatorname{HKDF\mbox{-}Expand}
(k_{\mathrm{root}},\texttt{assignment}\|m\|v\|R\|d), \\
k_{\mathrm{scope}} &= \operatorname{HKDF\mbox{-}Expand}
(k_{\mathrm{root}},\texttt{authorization}\|\mathcal{S}),
\end{align}
where the canonical scope tuple is
\begin{equation}
\begin{aligned}
\mathcal{S}=(&\text{issuer},\text{subject/sub-id},\text{tenant},
\text{regime/capability},m,\\
&\text{resource IDs},\text{operation},\text{grant ID},
\text{key epoch},\\
&\text{expiry},\text{policy version}).
\end{aligned}
\label{eq:canonical-scope}
\end{equation}

The assignment is generated once from $k_{\mathrm{assign}}$ and
stored by identifier.  A child capability does not independently
shuffle the model; it authorizes an assignment already present in
the global partition.  Tokens may encrypt that assignment identifier
and capability payload with AES-256-GCM while binding the canonical
scope as Associated Authenticated Data (AAD).  Tampering with a bound
field causes authentication failure.

Tenants do not choose or execute the Fisher--Yates shuffle.  The
lockbox controls derivation authority, and the trusted runtime
resolves an authenticated credential to a scoped capability and its
assigned mask.  A caller presenting a credential for regime $s$
cannot ask the runtime to apply regime $r$'s assignment.  Possessing
or guessing mask indices is not sufficient because the runtime does
not accept caller-supplied masks.

AAD authenticates claims; it does not evaluate them.  Immutable
fields such as tenant, sub-id, regime, model digest, operation, and
policy version are checked against the request during decryption.  A
grant may additionally authenticate administrative fields such as a
parent grant, allowed operations, an input/output-token ceiling, a
maximum invocation count, or a delegation right.  Dynamic
requirements such as a remaining-use counter, geography, current
account state, or revocation status must be evaluated by the policy
runtime using those authenticated claims.  In particular, placing
\texttt{max\_uses=N} in AAD makes the limit tamper-evident but does not
decrement it.  Exact-use enforcement requires an online atomic ledger
that reserves a use before invocation, rejects replay, and fails closed
when current state cannot be read.  The capability is therefore the
combination of an authenticated scope, lockbox grant, current policy
state, and enforcing runtime---not the mask vector alone.

\subsection{Gate Surfaces}
\label{sec:gate-surfaces}

The word \emph{gate} denotes two related enforcement surfaces, not one
claim stretched across every architecture.  Within an FFN, the gate is
a Hadamard mask on a declared activation.  Around a complete model, it
is a fail-closed execution decision.  The location determines the
state covered by the guarantee.

\paragraph{Pre-MLP gate (FFN input).}
Let $\bx\in\R^{d_{\mathrm{model}}}$ be the input to an MLP/FFN.  A
pre-MLP placement computes
\begin{equation}
\bar{\bx}=\bx\odot\bg_r,\qquad
\bh=\sigma(\bar{\bx}\bW_1+\mathbf b_1).
\label{eq:pre-mlp-gate}
\end{equation}
This is an FFN gate even though it precedes the first linear
projection.  Excluded input coordinates, their loss gradients, and
the loss gradients of the aligned input-facing rows of $\bW_1$ are
zero.  This placement alone does not partition the expanded hidden
activation, $\bW_2$, or parameters upstream of a mixing operation.

\paragraph{Intermediate-FFN gate.}
The stronger parameter-aligned placement computes
\begin{equation}
\bh=\sigma(\bx\bW_1+\mathbf b_1),\qquad
\tilde{\bh}=\bh\odot\bg_r,\qquad
\hat{\mathbf y}=\tilde{\bh}\bW_2+\mathbf b_2.
\label{eq:intermediate-ffn-gate}
\end{equation}
It zeros excluded hidden coordinates and aligns the loss-gradient
boundary with columns of $\bW_1$, entries of $\mathbf b_1$, and rows
of $\bW_2$.  The formal weight-confinement theorem and strict
DistilBERT experiment use this placement.

\paragraph{Whole-model gate.}
For model resource $M$ and operation $o$, let
$a(c,M,o)\in\{0,1\}$ be the trusted runtime's authorization decision
after resolving the authenticated tenant + sub-id scope.  The atomic
execution gate is
\begin{equation}
\operatorname{Exec}(c,M,\bx,o)=
\begin{cases}
M(\bx), & a(c,M,o)=1,\\
\bot, & a(c,M,o)=0,
\end{cases}
\label{eq:whole-model-gate}
\end{equation}
where $\bot$ means rejection before model invocation.  No internal
coordinate is removed on the authorized path, so this gate introduces
no model-capacity loss; it relies on the IdP, lockbox/proxy, runtime,
and serving-state boundary.  A model bank may contain many atomic
model resources, but that resource count is not the FFN partition
ratio $d/R$.  Whole-model authorization therefore does not claim
multi-regime separation inside one shared Transformer.  The decision
is nevertheless a structured capability rather than a bare Boolean:
it can bind the model version, operation, tenant + sub-id, expiry,
delegation chain, token ceiling, use budget, and audit identifier.
Several atomic gates may be chained around a pipeline---for example,
retrieval, reranking, generation, and a tool call---with each stage
requiring its own resource-and-operation grant.  The authorized model
computation remains unchanged at every stage.

\subsection{Hadamard Governance}
\label{sec:hadamard}

The FFN gate is governed by exactly one algebraic rule:

\begin{definition}[Hadamard gate]
\label{def:hadamard-gate}
For hidden activation $\bh \in \R^d$ and binary mask
$\bg_r \in \{0,1\}^d$:
\begin{equation}
\tilde{\bh} = \bh \odot \bg_r
\label{eq:hadamard-gate}
\end{equation}
\end{definition}

From this single definition, four governance properties follow,
each a direct consequence of $0 \cdot x = 0$ and $1 \cdot x = x$
propagated through the chain rule:

\begin{table}[ht]
\centering
\caption{The four governance properties of Hadamard gating.
Each is a one-line proof in Lean~4 (all reduce to \texttt{simp}).}
\label{tab:governance}
\begin{tabular}{@{}p{3.5cm}p{9cm}@{}}
\toprule
\textbf{Property} & \textbf{Statement} \\
\midrule
Forward isolation &
$g_{r,j} = 0 \implies (\bh \odot \bg_r)_j = 0$ \\
Active preservation &
$g_{r,j} = 1 \implies (\bh \odot \bg_r)_j = h_j$ \\
Gradient isolation &
$g_{r,j} = 0 \implies
(\partial \cL / \partial h_j) = 0$ \\
Gradient confinement &
$(\partial \cL / \partial h_j) \neq 0 \implies
g_{r,j} = 1$ \\
\bottomrule
\end{tabular}
\end{table}

The IdP-authenticated runtime decides which mask is authorized; the
Hadamard product enforces that resolved decision at the declared
tensor.  The runtime is therefore essential to authorization, while
$0\cdot x=0$ supplies the local coordinate exclusion and gradient
exclusion once the correct mask is applied.

\paragraph{Computational overhead of the gate.}
\label{sec:gate-overhead}
The gate is one element-wise multiply over $d$ coordinates per gated
tensor per invocation: $d$ scalar multiplications with no accumulation in the
forward pass and one broadcast multiply in the backward pass, against
the $\Theta(d\cdot d_{\text{in}})$ and $\Theta(d\cdot d_{\text{out}})$
MACs of the two projections it guards.  It is emphatically not an
extra matrix multiplication: implemented as a dense --- or even an
explicit diagonal --- matrix product $\mathrm{diag}(\bg_r)\,\bh$ it
would cost $\Theta(d^2)$ MACs per gated boundary, but that is an
implementation error, not a property of the mechanism.  Evaluated as
element-wise broadcast multiply (as in the checked-in implementations)
its memory traffic is two vector reads and one vector write and may be
fused into the epilogue or prologue of a surrounding GEMM.  For
inference-output equivalence, a mask that is constant for a packaged
artifact may instead be folded into the aligned rows of $\bW_2$.
That reparameterization removes the per-request multiply but does not
instantiate the declared activation-level zero or the training-state
confinement theorem; it must be described as extraction or packaging,
not as the same gate surface.  A bit-packed mask requires $d$ bits
(96 bytes at $d{=}768$), while an expanded Boolean or floating-point
runtime tensor requires implementation-dependent additional storage.
Enforcement is therefore not free, but the unfused operation is linear
in the gated width rather than quadratic.  We make no numerical latency claim for the
gate: the only systems figures in this repository are
\Cref{sec:deployment}'s state-dict arithmetic and
\Cref{sec:permutation-equivariance}'s keyed-basis benchmark, the
historical throughput figures are withheld pending hardware-matched
rerun (\Cref{sec:deployment}), and no dedicated gate-overhead
microbenchmark is yet part of the checked-in evidence.

\paragraph{Orthogonality of disjoint gated activations.}
For two regimes $r,s$ with disjoint masks
($\bg_r \odot \bg_s = \mathbf{0}$), the gated activation vectors
have zero inner product:

\begin{equation}
\langle \bh \odot \bg_r, \; \bh \odot \bg_s \rangle
= \sum_{j=1}^{d} h_j^2 \, g_{r,j} \, g_{s,j}
= 0
\label{eq:entanglement-breaking}
\end{equation}

because $g_{r,j}g_{s,j}=0$ for every coordinate.  This is an
algebraic support-orthogonality statement, not a claim that the
random variables represented by different regimes are statistically
independent or that shared model paths carry no correlated
information.

\subsection{Gated Forward Pass}
\label{sec:gated-forward}

\begin{definition}[Element-wise FFN independence]
\label{def:ffn-independence}
In a standard transformer FFN, each hidden dimension is computed
independently: $h_j = \sigma(\bx \cdot \bW_1[:,j] + b_{1,j})$.
The output contribution of dimension $j$ is
$\hat{y}_m = \sum_j \tilde{h}_j \, W_{2,j,m} + b_{2,m}$.
Zeroing $\tilde{h}_j$ via the gate eliminates dimension $j$'s
contribution to every output element exactly.
\end{definition}

For input $\bx$, regime $r$ with gate mask $\bg_r$:

\begin{align}
    \bh &= \sigma(\bx \bW_1 + \mathbf{b}_1)
    \label{eq:ffn-hidden} \\
    \tilde{\bh} &= \bh \odot \bg_r
    \label{eq:gate-apply} \\
    \hat{\mathbf{y}} &= \tilde{\bh} \, \bW_2 + \mathbf{b}_2
    \label{eq:ffn-output}
\end{align}

The output depends only on the $d_r=d/R$ dimensions
owned by regime $r$.  This is the architectural reason CDP
uses FFN dimensions as its core partition: they are independent,
whereas dimensions inside a dense shared attention head are coupled
through $\bQ\bK^\top$.  Attention isolation requires the separate
lane construction described in \Cref{sec:trust-boundary}, not a
post-hoc coordinate mask inside one head.

\subsection{Gradient and Weight Confinement}
\label{sec:gradient-isolation}

During backpropagation, the chain rule propagates the gate mask
into the gradient:

\begin{align}
    \frac{\partial \cL}{\partial \tilde{\bh}} &=
    \frac{\partial \cL}{\partial \hat{\mathbf{y}}} \bW_2^\top \\
    \frac{\partial \cL}{\partial \bh} &=
    \frac{\partial \cL}{\partial \tilde{\bh}} \odot \bg_r
    \label{eq:grad-isolation}
\end{align}

For any dimension $j \notin \mathrm{supp}(\bg_r)$:
$g_{r,j} = 0 \implies \partial \cL / \partial h_j = 0$.
This is exact, not approximate, and holds unconditionally
for all losses, batch sizes, and inputs at the gradient level.
Consequently:

\begin{itemize}[nosep]
    \item $\bW_1[:,j]$ receives zero loss gradient for
    $j \notin \mathrm{supp}(\bg_r)$ (proven:
    \texttt{weight\_update\_confined}).
    The corresponding hidden bias $b_{1,j}$ receives the same masked
    preactivation gradient.
    \item $\bW_2[j,:]$ receives zero loss gradient for
    $j \notin \mathrm{supp}(\bg_r)$ because $\tilde{h}_j = 0$.
    This is \texttt{w2\_update\_}\allowbreak\texttt{confined}.
\end{itemize}

The output bias $\mathbf{b}_2$ is not coordinate-indexed by the
hidden partition and is therefore shared unless frozen, separately
scoped per regime, or gated by an additional declared mechanism.  It
is not covered by the $\bW_1/\bW_2$ confinement theorem.

Exact parameter-state confinement additionally requires the
optimizer to respect the same coordinate boundary: decoupled
weight decay, stale momentum, and other stateful update terms must
be masked per coordinate or maintained per regime.  Subject to
that optimizer condition, updates to the gated FFN parameters are
confined to the regime's coordinate partition.  The
gate restricts both what the output can source (forward) and where
loss gradients can flow (backward); shared components outside this
surface are identified in \Cref{sec:isolation-surface}.

\subsection{Regime Destruction (Algebraic Annihilation)}
\label{sec:destruction}

Destroying regime $r$ consists of zeroing the weight columns and
rows corresponding to $\mathrm{supp}(\bg_r)$:

\begin{equation}
    \bW_1[:,j] \leftarrow \mathbf{0}, \quad
    b_{1,j} \leftarrow 0, \quad
    \bW_2[j,:] \leftarrow \mathbf{0} \quad
    \forall\, j \in \mathrm{supp}(\bg_r)
\end{equation}

By the confinement theorems, no regime-$r$ update exists outside
$\mathrm{supp}(\bg_r)$ within the confined FFN parameter surface.
After zeroing, those parameters are identical to parameters that
never received regime-$r$ loss gradients.  If $\mathbf{b}_2$, all
other tenant-trainable paths, and optimizer state are frozen,
separately partitioned, or erased under the same boundary, the statement
extends to the complete declared regime-scoped trainable state.  We
call this algebraic annihilation: mathematical zeroing of the
declared parameter slice.  It does not erase shared pretrained
representations, logs, caches, or backups.

\subsection{Cross-Mask Additive Cancellation}
\label{sec:additive-cancellation}

A narrow adversarial question is whether an additive vector confined
to one mask can appear through a disjoint mask at the same gate.  The
answer is no under the theorem below.

\begin{theorem}[Injection confinement]
\label{thm:injection-confinement}
Let $\mathbf{z} \in \R^d$ be a knowledge vector injected
additively into the residual stream, confined to regime $r$'s
coordinate subspace via the gate mask.  For any other regime
$s \neq r$, the gated output of regime $s$ is identical with
and without the injection:
\[
    (\bh + \mathbf{z} \odot \bg_r) \odot \bg_s = \bh \odot \bg_s
\]
Proven in Lean~4: \texttt{injection\_confined\_vec}.
\end{theorem}

The proof follows from two facts: (1)~$\mathbf{z} \odot \bg_r$
is zero at every dimension outside $\mathrm{supp}(\bg_r)$, and
(2)~$\mathrm{supp}(\bg_r) \cap \mathrm{supp}(\bg_s) = \emptyset$
for $r \neq s$.  The injection vanishes at every dimension regime
$s$ owns at that gate.  This does not cover information routed
through shared components or ungated bypasses.

Injection confinement complements aligned FFN gradient confinement:
one constrains an explicitly gated additive vector, while the other
constrains updates on the declared FFN parameter surface.  Neither
statement extends to shared or bypass paths.

\subsection{Mask Algebra and Authorized Union}
\label{sec:composition}

Binary masks are closed under AND, NOT, and OR.  This algebraic fact
does not make every operation an authorized runtime primitive:

\begin{definition}[Regime composition]
\label{def:composition}
For regimes $r, s$ with masks $\bg_r, \bg_s$:
\begin{align}
    \bg_{r \cap s} &= \bg_r \odot \bg_s
    \quad\text{(AND: element-wise product)}
    \label{eq:and} \\
    \overline{\bg}_r &= \mathbf{1} - \bg_r
    \quad\text{(NOT: element-wise complement)}
    \label{eq:not} \\
    \bg_{r \cup s} &= \overline{\,\overline{\bg}_r \odot \overline{\bg}_s\,}
    \quad\text{(Union: derived via De~Morgan)}
    \label{eq:union}
\end{align}
AND and NOT are sufficient for functional completeness; Union is a
convenience derived as $\bg_{r\cup s}=\mathbf{1}-(\mathbf{1}-\bg_r)\odot(\mathbf{1}-\bg_s)$.
For disjoint regimes, AND returns the zero vector and Union
returns the concatenation of both supports.
\end{definition}

The production-facing primitive considered here is an explicit
\textbf{authorized union}, which activates both supports in one forward
pass.  The Lean~4 declarations \texttt{compose\_disjoint} and
\texttt{compose\_access} establish the mask-level facts that each
constituent support remains active and coordinates outside the union
remain excluded.  They do not prove end-to-end safety of shared model
components.

The complement $\overline{\bg}_r$ activates every coordinate outside
$r$ and therefore cannot be inferred as an authorization grant from
possession of $r$ alone.  A runtime may construct such a mask only if
the resolved scope independently authorizes every activated support.

Union is an authorized operation: the caller must possess a valid
grant for every constituent regime.  Those grants may belong to
different tenants who consent to composition, or to one principal
that is intentionally assigned several capabilities (for example,
multiple MoE experts).  Under the trusted-runtime assumptions, the
runtime refuses to construct the Union mask without all required
grants.  A multi-regime grant deliberately removes the FFN-coordinate
boundary among its members while preserving mask-level exclusion for
coordinates outside the union.

\paragraph{Whole-model wrapping.}
A complete model may be registered as one atomic resource behind
\eqref{eq:whole-model-gate}, regardless of how many other models the
runtime serves.  This is not a degenerate FFN mode: the enforcement
event is authorization before invocation, not multiplication of every
internal activation by an all-ones mask.  If an inner FFN partition is
also configured with $R_{\mathrm{FFN}}{=}1$, its mask
$\bg_0=\mathbf 1$ is a computational no-op, but the outer model gate
still supplies identity, operation, lifecycle, delegation, metering,
and audit control.  It can also serve as one link in a capability
chain: authorization to invoke model $M_1$ need not imply authorization
to invoke a downstream model $M_2$, retrieve corpus $D$, or call tool
$T$.  Each edge is resolved from its own AAD-bound resource and
operation scope.
Later internal partitioning requires separate gate-aware validation
and may change utility; atomic whole-model authorization does not
predict $R_{\mathrm{FFN}}>1$ behavior.

\paragraph{Nested composition (CDP inside CDP).}
An outer whole-model gate and an inner FFN gate compose sequentially:
the runtime first authorizes model $M$, then $M$ executes with the
authorized inner mask $\bg_r$.  The outer decision acts on a resource
identifier; the inner mask acts on an activation tensor.  They are not
one Hadamard vector, and the Lean~4 coordinate proofs cover the inner
FFN gate rather than the outer control-plane decision.

When multiple masks are defined over the same indexed activation space,
associativity permits several masks to reduce to one binary vector.
If outer and inner policies act on different
axes---for example expert identity and coordinates inside an
expert---they must first be lifted into one product space or a
block-diagonal mask.  Their original vectors cannot be multiplied
directly.

This nested structure is particularly natural for
Mixture-of-Experts (MoE) architectures, where each expert is
already a separate FFN block.  A CDP partition at the expert level
can provide \emph{capability selection} (which experts a principal can
activate), while a CDP partition within each expert provides
\emph{coordinate confinement} (which dimensions inside an expert are
active).  Under the product-space lifting above, the two levels
compose because disjointness is preserved under Hadamard
product: if $\bg^{\mathrm{outer}}_i \odot \bg^{\mathrm{outer}}_j
= \mathbf{0}$ and $\bg^{\mathrm{inner}}_r \odot
\bg^{\mathrm{inner}}_s = \mathbf{0}$, then the composed masks are
also disjoint.

\paragraph{The expert-routing plane.}
\label{sec:moe-routing}
Between the two levels sits an expert-router authorization surface.
A MoE layer selects experts through a routing projection that produces
expert logits and dispatch weights.  Enforcement must respect that
router's order of operations.  One conforming construction first
restricts candidates to the authorized set $A$ (equivalently, sets
every unauthorized logit to $-\infty$), selects at most
$k'=\min(k,|A|)$ finite entries, and refuses to dispatch any
$-\infty$ entry.  Another multiplies normalized dispatch weights by a
grant-scoped binary mask before dispatch and removes every zero-weight
entry.  Both must fail closed when $A$ is empty.  Plain multiplication
of pre-softmax logits by zero is insufficient because softmax can
assign positive mass to a zero logit; setting logits to $-\infty$
without constraining a fixed-size top-$k$ is likewise insufficient
when $|A|<k$.  Exact zero contribution follows algebraically when
normalized dispatch weights are masked; exclusion from selection and
computation additionally depends on the fail-closed router
implementation.  The distinction to keep straight is
one of control versus routing: MoE routing is data-dependent and
learned, deciding \emph{which tokens} see which experts; CDP
authorization is data-independent and grant-scoped, deciding
\emph{which principals} may reach which experts at all.

We evaluate this composition on eight disjoint binary tasks from the standard
20 Newsgroups train/test split. Each regime independently trains a top-1 router
over two private feed-forward experts with task loss and a load-balancing loss
computed only inside its authorized set. The eight trained routers and sixteen
experts are copied unchanged into one execution-authorized bank. Across 6,231
held-out examples, macro accuracy is 81.80\% and micro accuracy is 81.86\%
(5,101 correct). Separate and packed execution have exactly equal logits and
predictions; unauthorized expert dispatch is zero. Both experts are used in
every regime, with the least-used expert receiving 42.4\% of its regime's
held-out examples, and every router records nonzero parameter movement. A
dedicated active-regime training-step audit records zero inactive gradient
tensors, optimizer states, and parameter change. Wrong regime, model,
dimension, and empty-authority probes against eight single-regime execution
surfaces make zero model calls.

The deliberately invalid control multiplies unauthorized pre-softmax logits by
zero; an unauthorized expert wins top-1 because zero exceeds the authorized
negative logits. Candidate restriction or $-\infty$ masking before softmax is
therefore an empirically exercised placement requirement. Exact zero dip here
is a copy-and-pack equivalence result. The packed bank has the same
router/expert bytes as the separate lanes, so this experiment does not itself
claim storage compression, shared-attention isolation, or language-model
quality.

\begin{sloppypar}
\noindent\textit{Replay learned routing.} Run
\path{experiments/train_authorized_moe.py}; the clean-source artifact is
\path{experiments/results/authorized_learned_moe_20260831T182431_055309Z.json}.
\end{sloppypar}

\paragraph{Execution-selected capability prefixes and token routing.}
Learned continuous prefixes can condition model computation without changing
the user-visible token sequence~\citep{li2021prefix}. They are not, by
themselves, credentials. Any plaintext prefix accepted from a request is
forgeable and remains untrusted data. Our construction therefore separates the
channels: authenticated execution authority selects a private learned prefix and
its allowed expert set; user tokens then condition top-1 routing only within
that set. The prefix enters the token router as a private continuous context,
not as user-authored text, and authorization is applied before softmax and
dispatch.

We evaluate eight private prefixes and sixteen experts on the same held-out
20 Newsgroups task family. Across 311,477 non-padding token decisions, macro
accuracy is 77.76\% and micro accuracy is 77.80\% (4,848 of 6,231 examples).
Copying the eight independently trained token routers, prefixes, and experts
into one authorized bank changes zero token routes, logits, or predictions.
Global negative-infinity masking exactly matches candidate restriction;
unauthorized dispatch is zero. A spoof request containing literal
\texttt{CAPABILITY\_REGIME\_0} through
\texttt{CAPABILITY\_REGIME\_7} as user text also causes zero unauthorized
dispatch. Every learned prefix and router moves, both experts are used in every
regime with a minimum held-out share of 41.7\%, and an active-regime update
creates zero inactive gradients, optimizer states, or parameter change. All
malformed probes against eight singleton execution authorities make zero model
calls.

This is a small token-routing classifier, not an autoregressive Transformer.
The result does not establish prompt text as authority, language-model
generation quality, shared-attention isolation, or storage compression. Exact
zero dip is again the copy-and-pack equivalence; task utility is reported
separately.

\begin{sloppypar}
\noindent\textit{Replay capability-prefix routing.} Run
\path{experiments/train_capability_token_moe.py}; the clean-source artifact is
\path{experiments/results/capability_prefix_token_moe_20260831T182453_013039Z.json}.
\end{sloppypar}

% ===================================================================
\section{Threat Model and Trusted Execution}
\label{sec:threat-model}

\subsection{System and Assets}

The system input is not a prompt $\bx$ alone but a pair $(\bx,c)$,
where $c$ is an authenticated user or workload credential.  The IdP
and issued lockbox grant resolve $c$ to the canonical scope $\mathcal S(c)$
from \eqref{eq:canonical-scope}; tenant identity without the
authenticated subject/sub-id is incomplete.  A
principal may hold one regime or an authorized set of capabilities
$S\subseteq\{0,\ldots,R-1\}$.  The trusted control path resolves
\begin{equation}
c \longmapsto
\mathcal S(c),\quad S,\quad \{(k_{\mathrm{scope},r},\bg_r):r\in S\},
\quad \bg_S=\bigvee_{r\in S}\bg_r
\label{eq:credential-resolution}
\end{equation}
before model execution, where $\bigvee$ is element-wise binary OR.
For a single-regime credential, $S=\{r\}$ and $\bg_S=\bg_r$.
For a whole-model gate, the same canonical scope instead resolves a
model resource $M$ and operation $o$ and requires $a(c,M,o)=1$ before
invocation; no FFN mask is implied by that authorization.
The union is defined only for regime masks expressed in the same
active model basis.  The residual-stream permutations $P_r$
introduced in \Cref{sec:orthogonal-superposition} are alternative
model-wide keyed placements,
not coordinates that may be ORed directly.  A grant spanning
different keyed basis placements must execute each placement
separately and compose only at a declared output boundary; the current
runtime does not claim single-pass union across distinct $P_r$ bases.
The protected assets are the root and assignment keys, lockbox
authority, scoped capability keys, permutation assignments, masks,
declared regime-scoped FFN coordinates, model attestations, and the
audit record binding those objects to the complete canonical scope.

The intended trust chain is
\[
\text{IdP} \to
\text{lockbox/Vault} \to \text{sidecar} \to
\text{runtime} \to \text{gate}.
\]
The implementation is not coupled to a particular trusted-compute
vendor.  Its portable machine-credential path uses PKCS\#12 to package an
Ed25519 private key, leaf certificate, and certificate chain, while the
verifier pins the accepted root independently.  PKCS\#12 is a credential
container rather than evidence of TPM, enclave, measured-boot, or
confidential-computing residency.  The reference path loads the key into
process memory; deployments requiring non-exportable keys must supply a
native signing provider and separate platform evidence.  A deliberately
pinned self-signed root is also valid, but establishes only possession and
continuity of that pinned key.

This chain contains two identities that must not be conflated.  An
OIDC IdP authenticates the human or calling principal and supplies the
issuer and stable subject/sub-id.  SPIFFE/SPIRE supplies short-lived,
attested workload identities (SVIDs) and trust bundles to the
lockbox, sidecar, and runtime workloads~\citep{spiffe2026}.  SPIFFE is
therefore the recommended service-to-service authentication layer; it
does not replace the end-user IdP or decide which tenant resource the
user may access.

The IdP-authenticated principal resolves to a tenant, resource, and
operation scope at the lockbox or Vault-backed provider, which controls
principal-scoped key derivation and grant issuance for the authorized
regime or capability set.  A sidecar must authenticate both the
current grant and the workload allowed to redeem it, release only the
authorized mask or invocation decision, and never accept a
caller-selected permutation.  The runtime is the leaf of this chain:
it applies the resolved gate immediately at the declared surface and
namespaces serving state under the same canonical scope.

For a metered whole-model grant, the strict request order is:
validate the OIDC principal; authenticate the enforcing workload with
an SVID; atomically reserve the grant's invocation and
token budget; obtain the current mask or model-execution decision;
invoke the declared resource; settle the measured token count; and
append an audit event.  SPIFFE authenticates the workload performing
those steps.  The atomic ledger, rather than the SVID or AAD alone,
enforces an exact number of uses.

\subsection{Adversary and Trust Assumptions}

The in-scope adversary is an unauthorized tenant or network caller
who can submit chosen inputs, observe released outputs, train on
data within its own authorized regime set, and possess its own scoped
credential or capability material.  The adversary may substitute a
tenant, sub-id, regime, model, resource, or operation field; tamper
with a token or model volume; or
attempt to compose regimes without the required grants.  It does
not control the IdP, lockbox, Vault, sidecar, trusted runtime
process, or master gate seed.

Compromise of that trusted custody boundary, process-memory
extraction after a mask is loaded, or an operator intentionally
issuing an unauthorized grant is outside the cryptographic threat
model.  It is nevertheless an explicit system risk.  The term
\emph{lockbox} covers both a sealed distributable artifact and the
authority that holds the root, assignment, and signing material.
Stealing a sealed artifact alone does not reveal recipient-wrapped
keys.  Exfiltrating the live authority, root/assignment key, or an
unrestricted derivation-and-signing capability is a keys-to-the-kingdom
failure: the attacker can reconstruct assignments and mint apparently
valid grants for every scope.  Recovery requires disabling the old
authority, rotating the root and accepted key epoch, reissuing grants,
and, where assignment secrecy matters, regenerating and reattesting the
partitioned model.

Key assignment is therefore an operator-controlled part of the
trusted computing base: isolated scopes require distinct capability
contexts, while key reuse or a multi-regime grant intentionally joins
their capability boundary.  Likewise, components declared shared in
\Cref{tab:isolation-surface} are not silently promoted into the
algebraic guarantee: KV-cache blocks, optimizer state, logs, and
other stateful artifacts must be separately namespaced by the same
resolved regime identity.  Gate placement and path closure remain
mandatory as described in \Cref{sec:isolation-surface}.

\subsection{Credential Binding and Fisher--Yates Custody}

Callers never choose the Fisher--Yates permutation.  The lockbox
controls the global assignment and capability authority, only scoped
material is released, and the runtime obtains the mask through the
authenticated sidecar path.  In the strict path, the OIDC validator
resolves issuer and subject/sub-id, the grant binds tenant,
regime/capability, model digest, resource, operation, grant ID, key
epoch, and expiry, and SPIFFE authenticates the
workload that requests or redeems the grant.  The complete scope must
match the request, the current grant, and the stored artifact before
mask indices or an invocation decision are returned.  A
caller-supplied \texttt{regime\_id} is not sufficient.
Multi-capability execution requires an authenticated grant
covering every member of $S$; the current single-operator path uses
an administrator-authorized compound request, while strict
cross-authority grants remain an integration seam.  The sidecar
separately verifies ownership before proxying a request.  Vault ACLs and lockbox
grants govern operator and workload access to derivation operations.

Deriving a separate key for each principal makes revocation granular,
but derivation does not claw back bytes already disclosed.  A strict
grant therefore includes a unique grant ID and monotonically changing
key epoch in its HKDF context.  Revoking one principal marks that grant
inactive and advances its accepted epoch; every mediated invocation
checks current state before use.  Other principals need not rotate.
Removing a recipient from a newly issued lockbox, by itself, does not
invalidate an old derived key copied for offline use.  A deployment
that accepts such keys without an online current-epoch check cannot
claim immediate principal revocation.

Therefore guessing a permutation is not itself an authorization
event.  Dispatch integrity relies on the IdP,
lockbox, sidecar, and runtime boundary; assignment secrecy under the
HMAC-SHA256 PRF assumption is defense in depth.  The public
reproduction bundle contains the FFN gates, a minimal fail-closed Cargo
scope harness, and an atomic whole-model invocation check.  Production
OIDC, Vault, and sidecar
implementations are outside this repository and are not claimed as
independently reproduced here.

\subsection{Guarantees Under Correct Dispatch}

For a whole-model scope, correct dispatch either invokes the complete,
unchanged model or rejects before invocation.  This is an operational
authorization guarantee under the trusted runtime boundary, not an
internal coordinate-separation theorem.

For an invocation resolved to capability set $S$ in one active model
basis, the runtime applies $\bg_S$.  For every coordinate owned by a
regime $r\notin S$,
disjointness makes the corresponding activation exactly zero at
the correctly placed gate.  The full activation need not be the
zero vector: coordinates in the authorized union remain active.
Thus the structural claim is non-access outside $S$.  A
single-capability invocation commonly uses $|S|=1$; a capability or MoE
deployment may intentionally use $|S|>1$.  In the latter case,
sharing inside $S$ is authorized behavior, not cross-tenant leakage.
Any wrong-key task-accuracy effect is downstream and empirical rather
than part of the gate theorem.

Matched conjugation and wrong-regime execution are also distinct
conditions.  A complete conjugation $P_r$ paired with its matching
runtime basis preserves the base function exactly.  A mismatched
credential, basis, mask, head, or classifier leaves an uncancelled
coordinate relabeling and does not reproduce the intended regime.
Permutation conjugation provides the keyed placement; the Hadamard
gate provides zeros on excluded FFN coordinates.

The operator is responsible for preserving the intended key
cardinality.  If two regimes that are claimed to be isolated are
assigned the same effective key, mask, or unrestricted compound
grant, CDP correctly enforces the resulting joined capability but
cannot infer that the policy was mistaken.  The runtime reduces this
risk by deriving keys in regime-specific contexts, validating
identity-bound grants, and refusing caller-selected permutations.

\subsection{Prospective Knowledge Lifecycle and BYOK}

A provider could reserve an unassigned FFN partition for later
regime-scoped training.  The confinement theorem predicts that a
correctly placed intermediate FFN gate and coordinate-respecting
optimizer would restrict parameter updates to that partition.
However, the present experiments do not establish successful
owning-key learning in a previously unused partition.  Hollow-regime
installation, tenant-corpus deposition, per-regime heads or adapters,
and injected encoders are therefore prospective embodiments rather
than evaluated contributions of this paper.

In the BYOK embodiment, the tenant additionally supplies an inner
key that derives $\bg^{\mathrm{tenant}}$.  The effective mask is
\[
\bg^{\mathrm{effective}}_r =
\bg^{\mathrm{regime}}_r \odot \bg^{\mathrm{tenant}}.
\]
The platform can authorize the outer regime without possessing the
tenant's inner key.  At the gate, a coordinate is active only when
both masks include it.  Tenant-specific content utility under that
construction is prospective.  Nested-mask and portable PKCS\#12 key
provider paths exist as reference implementations; strict hardware-backed
key custody and
end-to-end BYOK integration with every production runtime path is an
embodiment boundary, not an assumption of the core gate theorem.

\subsection{Mandatory Enforcement Architecture}

CDP is an architecture comprising the authenticated IdP
route, lockbox or Vault-backed derivation, enforcing sidecar, trusted
runtime, and correctly placed gate.  A conforming deployment must
resolve the authenticated canonical scope through that chain; an
unverified request \texttt{tenant\_id}, \texttt{sub\_id}, or
\texttt{regime\_id} is not an authorization input.
Development routes that accept direct identifiers or injected keys
bypass CDP and cannot claim its enforcement properties.

\paragraph{Audited implementation status.}
The public release contains the Gate library and no serving or control-plane
implementation. It provides deterministic mask derivation, scoped AEAD and
HMAC contracts, recipient-wrapped grants, signed lockboxes, exact signed-grant
membership checks, PKCS\#12/X.509 verification against operator-configured
roots, finite operation gates, and the fail-closed Cargo protocol. Cargo binds
tenant, subject, Regime, model, operation, policy version, payload semantics,
and exact partition; its success receipts require declared load completion and
commit to material retrieval results.

Current experiments use a small source-visible callback preflight that checks
model identity, dimensions, and Regime before invoking a model. This fixture is
not an IdP, credential verifier, network service, durable use ledger, or
production authorization layer. Accordingly, the repository makes no claim
that an end-to-end production principal path is shipped.

The production profile requires an operator to authenticate the principal,
resolve a signed canonical scope, distribute trusted roots, protect keys,
enforce expiry and revocation, redeem finite uses atomically across replicas,
apply the Gate before every protected operation or declared FFN path, and close
all alternate routes. Those deployment properties require separately pinned
evidence from the actual integration.

% ===================================================================
\section{The Trust Boundary}
\label{sec:trust-boundary}

\subsection{Isolation Surface by Embodiment}
\label{sec:isolation-surface}

CDP's \emph{isolation surface} is the set of activations,
parameters, optimizer state, and serving state for which a
deployment enforces a regime boundary.  The surface is
embodiment-specific; the binary gate always gives exact forward
and loss-gradient confinement at the activation on which it acts.
That local algebra does not certify an arbitrary installation point.
For the declared surface to be isolated, the gate must precede the
first operation that mixes protected coordinates into unprotected
ones, and every path carrying the protected representation must
cross an equivalent gate.  A downstream or bypassed gate can still
zero its immediate tensor while failing to isolate the computation
as a whole.

\begin{table}[ht]
\centering
\caption{Isolation surface by gate placement.  Pre-MLP and
intermediate placements are both FFN gates but cover different
parameter axes.  Whole-model gating is an operational invocation
boundary, not an internal coordinate partition.}
\label{tab:isolation-surface}
\scriptsize
\begin{tabularx}{\textwidth}{@{}p{3.0cm}XXX@{}}
\toprule
\textbf{Component or state} &
\textbf{Pre-MLP FFN} &
\textbf{Intermediate FFN} &
\textbf{Whole model, atomic} \\
\midrule
Gated tensor &
FFN input coordinates &
Expanded hidden coordinates &
No internal tensor mask \\
Aligned loss-gradient and parameter surface &
Input gradient and $\bW_1$ input-facing rows &
$\bW_1$ hidden columns, $\mathbf b_1$, and $\bW_2$ rows &
Complete model invoked or rejected as one resource \\
Attention, embeddings, normalization, output paths &
Shared; outside local gate &
Shared; outside local gate &
Inside authorized atomic resource; not internally partitioned \\
Optimizer moments and decoupled weight decay &
Mask aligned $\bW_1$ rows &
Mask aligned FFN slices &
Scope complete model-training state if training is authorized \\
KV cache, adapters, extracted artifacts, logs &
Execution-scoped &
Execution-scoped &
Model-scope and principal-scope runtime namespacing \\
Authorized-path capacity effect &
Uses active FFN-input support &
Uses active expanded-FFN support &
None: authorized model computation is unchanged \\
\bottomrule
\end{tabularx}
\end{table}

The local learned-state claim follows the declared FFN placement.  A
pre-MLP gate covers input-facing $\bW_1$ rows; the stronger evaluated
intermediate gate covers aligned $\bW_1$, $\mathbf b_1$, and $\bW_2$
slices.  Shared pretrained infrastructure may supply common
computation to all regimes.  If shared components are co-trained on
regime data, their updates are outside either local FFN guarantee.
Successful content-specific deposition is not established by the
current experiments.  Stateful
artifacts such as optimizer moments and KV-cache blocks are not
partitioned by multiplication alone and must follow the same
runtime-resolved regime identity.

\paragraph{Whole-model authorization.}
The entire model---attention, embeddings, FFNs, output projection,
normalization, and associated serving state---is registered as one
atomic resource.  Isolation is operational: the key hierarchy,
IdP-bound policy, runtime, audit, state namespacing, and lifecycle
controls must reject unauthorized invocation and state access.  A
runtime may serve many such model resources, so this mode is not
defined by $R_{\mathrm{FFN}}{=}1$.  Nor does it separate two
principals who are intentionally authorized to the same atomic model.
Its value is the structured authorization envelope around unchanged
computation: tenant + sub-id, model/version, operation, chain parent,
expiry, token ceiling, use budget, and audit identity can all be
bound without reducing the model's hidden dimension.  Exact cumulative
use or token limits additionally require the online reservation ledger
described above.

\paragraph{Intra-model FFN partitioning ($R_{\mathrm{FFN}}>1$).}
When multiple regimes share a single backbone, the gate acts at
FFN boundaries, either before the MLP or at its expanded intermediate
activation.  The following components remain shared across
regimes and are \emph{not} covered by the per-regime confinement
guarantee:

\begin{itemize}[nosep]
    \item \textbf{Attention mechanism.}  $\bQ\bK^\top$ computes a
    full inner product; the gate cannot act before the entanglement.
    \item \textbf{Token embeddings.}  The input embedding table is
    shared; all regimes map the same vocabulary to the same initial
    vectors.
    \item \textbf{Output projection.}  The language model head
    (unembedding matrix) is shared.
    \item \textbf{Layer normalization.}  LayerNorm statistics are
    computed over all active dimensions.
\end{itemize}

For the intermediate placement, the theorem confines the aligned FFN
fine-tuning delta when shared components are frozen and the optimizer
respects the coordinate partition.  The pre-MLP placement has the
narrower surface in \Cref{tab:isolation-surface}.  The present
classification experiments do not establish content-specific
deposition.  Shared components remain infrastructure inside the
trust boundary.

The honest summary is surface-specific.  A whole-model gate provides
atomic authorization without changing the authorized computation.  A
pre-MLP gate zeros selected FFN-input coordinates.  The evaluated
intermediate gate proves structural coordinate exclusion and aligned
FFN-slice confinement, while the shared backbone and runtime-scoped
state remain inside the declared trust boundary.

\paragraph{Attention as a duplicated lane.}
The evaluated construction leaves attention shared and gates only
the FFN intermediate activation.  In principle, attention could be
moved outside the shared boundary only by duplicating complete
Transformer lanes, including normalization and every residual route,
and gating before recombination.  That architecture has substantial
compute and operator-error risk and is outside this paper's evaluated
scope.  No result below relies on multi-regime attention isolation.

\paragraph{Gate placement constraint.}
A gate's location is part of the security mechanism, not a
performance-neutral implementation choice.  If protected signal is
mixed into active coordinates before masking, the later zero cannot
undo that transfer; if a residual, adapter, cache, or other branch
bypasses the gate, that branch lies outside the isolation surface.
Improper placement can therefore nullify the claimed isolation even
though multiplication at the misplaced gate still produces literal
zeros.

A concrete negative result sharpens this constraint.  Placing LayerNorm
\emph{between} the activation $\bh$ and the gate
$\tilde{\bh} = \bh \odot \bg_r$ breaks $\bW_1$ confinement
(proven: \texttt{non\_preserving\_breaks\_}\allowbreak
\texttt{w1\_confinement} in
Lean~4).  LayerNorm normalizes across all dimensions, mixing
information from gated-off dimensions back into gated-on
dimensions before the gate can act.  The safe placement is
$\text{act} \to \text{gate} \to \bW_2$, with LayerNorm either
before $\bW_1$ or after $\bW_2$.  For the narrower pre-MLP
embodiment, the gate must be immediately upstream of $\bW_1$ with no
intervening mixing operation; it does not inherit the intermediate
$\bW_1/\bW_2$ slice theorem.

% ===================================================================
\section{Formal Verification in Lean~4}
\label{sec:formal}

The checked-in proof inventory contains 21 Lean modules and 312 active
\texttt{theorem}/\texttt{lemma} declarations.  A clean build under
the pinned Lean/Mathlib toolchain contains no active \texttt{sorry}
or \texttt{admit} declarations.  The core Hadamard-gate, gradient,
and aligned FFN weight-confinement theorems are unconditional
algebraic results; a checked-in \texttt{\#print axioms} audit shows
that the nine headline algebraic and placement theorems use only
Lean/Mathlib's standard logical axioms.

The source tree separately contains five project-specific
\texttt{axiom} declarations.  Two declare opaque predicates and one
states a conditional PRF enumeration bound:
\begin{itemize}[nosep]
    \item \texttt{Recovers} (opaque predicate)
    \item \texttt{prf\_brute\_force\_optimal}
    \item \texttt{IsDistributionMatched} (opaque predicate)
\end{itemize}
They are not premises of the core gate-zeroing theorems.

Two further statistical axioms are retained in historical
conditional modules so that the old theorem chain remains
inspectable:
\begin{itemize}[nosep]
    \item \path{camouflage_indistinguishable}
    \item \path{gradient_probing_hard}
\end{itemize}
They are excluded from this paper's claim set and from the headline
axiom audit. Their justifications required
statistical distribution-matching of injected noise, which is not
provable against white-box weight inspection: an adversary holding
the model locally can classify dimensions by \emph{function} (coherent
activations and gradient signal on regime-appropriate inputs), and
marginal distribution matching does not constrain functional
distinguishers. Historical custody proposals are not part of the bundled
theorem inventory, so this release makes no custody theorem claim. Weight
opacity arising from cotraining is retained only as an \emph{observed}
hardening property---an empirical claim measured, not a theorem proved---and
is a basis for no security statement in this paper.

\subsection{Proof Architecture}

Selected proof modules are summarized below:

\begin{table}[ht]
\centering
\caption{Lean~4 proof files and their scope.}
\label{tab:proof-files}
\begin{tabular}{@{}p{4.5cm}p{8.5cm}@{}}
\toprule
\textbf{File} & \textbf{Scope} \\
\midrule
\texttt{GateSecurity.lean} &
Core isolation: forward/gradient isolation, weight confinement,
mask orthogonality, wrong-mask-reads-zero \\
\texttt{GatePlacement.lean} &
Placement constraints: support preservation, LayerNorm negative
result, residual block compositionality \\
\texttt{ModelSecurityV3.lean} &
Key hierarchy: confinement composition and recovery-conditioned
chain under the PRF axiom; earlier weight- and
partition-indistinguishability claims are excluded from this paper's claim
set, with their conditional declarations retained for audit (see axiom list) \\
\texttt{CapacitySecurity.lean} &
Cross-mask additive cancellation under disjoint support \\
\texttt{BiometricSecurity.lean} &
Biometric enrollment certificate and gated voice-separation conditions \\
\texttt{MultiEncoder.lean} &
Multi-encoder: federation, Procrustes alignment \\
\texttt{AttentionLeakage.lean} &
Negative result: coordinate masks do not isolate a dense shared
attention head \\
\bottomrule
\end{tabular}
\end{table}

\subsection{Selected Key Theorems}

\begin{table}[ht]
\centering
\caption{Selected theorems from the Lean~4 proof suite.}
\label{tab:lean-theorems}
\small
\begin{tabular}{@{}p{5.2cm}p{7.8cm}@{}}
\toprule
\textbf{Theorem} & \textbf{Statement} \\
\midrule
\texttt{gradient\_isolation} &
If $g_{r,j} = 0$, then the gradient w.r.t.\ $h_j$ is
exactly zero, unconditionally. \\
\texttt{weight\_update\_confined} &
Weight updates to $\bW_1$ are confined to columns in
$\mathrm{supp}(\bg_r)$. \\
\texttt{w2\_update\_confined} &
Weight updates to $\bW_2$ are confined to rows in
$\mathrm{supp}(\bg_r)$. \\
\texttt{masks\_orthogonal} &
For $r \neq s$: $\langle \bg_r, \bg_s \rangle = 0$. \\
\texttt{access\_requires\_exact\_mask} &
If a candidate mask reproduces the gated output for every activation
and every $\bW_2$, it equals the assigned mask coordinate-wise. \\
\texttt{compose\_disjoint} &
Union of disjoint masks preserves constituent access. \\
\texttt{residual\_block\_full\_}\allowbreak\texttt{confinement} &
FFN residual-branch backward chain confines aligned
$\bW_1/\bW_2$ updates for arbitrary upstream gradient. \\
\texttt{non\_preserving\_breaks\_}\allowbreak
\texttt{w1\_confinement} &
LayerNorm between activation and gate breaks confinement. \\
\texttt{adversary\_sees\_zero} &
For $r \neq s$, applying $\bg_s$ gives exactly zero on every
coordinate owned by regime $r$; the active coordinates are those
owned by $s$. \\
\texttt{rejection\_sampling\_count} &
Residue classes in the accepted rejection-sampling range have equal
cardinality; PRF pseudorandomness is a separate assumption. \\
\bottomrule
\end{tabular}
\end{table}

\subsection{Refinement Map and Verification Boundary}

The Lean vectors and matrices are linked to the strict DistilBERT
runner as follows.  With PyTorch's output-by-input linear-weight
layout, mathematical $\bW_1[:,j]$ corresponds to row $j$ of
\texttt{ffn.lin1.weight}; mathematical $\bW_2[j,:]$ corresponds to
column $j$ of \texttt{ffn.lin2.weight}.  The activation $\bh$ is the
tensor presented to each \texttt{ffn.lin2} forward pre-hook, and
$\bg_r$ is the Boolean 3,072-coordinate mask cast to that tensor's
dtype.  The research \texttt{ScopedAdam} stores first and second
moments at the same authorized entries.  Execution scope resolution
selects the mask before the model call.

Lean verifies the algebraic specification, not the Python hooks,
tensor-axis mapping, optimizer code, IdP, or absence of undeclared
bypasses.  The executable experiments therefore snapshot every
frozen and inactive tensor and assert exact zero deltas.  A conforming
implementation should additionally require deliberately misplaced-gate
and bypass controls to fail.  The public CPU harness includes both a
LayerNorm-before-gate control and an ungated-residual control.  These
controls exercise the Python placement logic; the Lean LayerNorm
counterexample proves the corresponding mathematical failure, but
neither is a machine-code verification of a deployment.

The key security assumption underlying the key hierarchy is
that HMAC-SHA256 is a pseudorandom function.

% ===================================================================
\section{Support-Constrained Training}
\label{sec:superposition}

An unexpected empirical finding is that the gate's measured
performance cost is much smaller than a naive capacity argument
would predict.  Halving the available dimensions might, in a linear
model, appear to halve representational capacity.  In five-seed
paired runs, co-trained DistilBERT remains within 0.14--0.38
percentage points of separately trained models across
$R{=}8$--128, including $R{=}96$ with only 8 dimensions per regime
on a 768-dimensional model.  Each independently trained inscription
has ordinary optimization variance; the empirical question is the
size and uncertainty of the performance cost, not exact equality.

We call this \emph{support-constrained training}.  Let
$A_r=\mathrm{supp}(\bg_r)$.  The gated FFN
or classifier contribution at a gated vector is
\[
  \sum_{j\in A_r} h_j\,\bW[j,:],
\]
because every term with $j\notin A_r$ is multiplied by zero.
Therefore every loss-reducing FFN contribution available to regime
$r$ must be represented through its active support $A_r$.  Across
jointly optimized regimes, each disjoint support must be sufficient
for its assigned task examples.  This is a direct linear-algebraic
consequence of the gate, not a causal hypothesis about a particular
microscopic feature geometry.  The empirical question is whether
optimization can satisfy the constraint at acceptable task cost;
\Cref{sec:series1} shows that it does at the formative post-encoder
classification surface.  \Cref{sec:strict-cotenancy} then instantiates
the constraint at every intermediate FFN of a frozen shared
DistilBERT backbone.  Its monotone accuracy decline from $R{=}1$ to
$R{=}16$ directly measures the cost of reducing each tenant from
3,072 to 192 FFN coordinates per layer.

\subsection{Gate-Aware Task Training}

Two training protocols are evaluated.  The formative five-seed models
are optimized for the downstream classification task with a mask on
the final pooled representation during every forward pass.  Because
that mask is downstream of the encoder, the protocol does not test
intermediate-FFN gradient or parameter confinement.  The strict
three-seed protocol instead freezes every non-governed shared path,
places a mask before each FFN down projection, updates only aligned
FFN parameter slices with regime-scoped optimizer state, and selects
a private classifier under the same regime identity.

% ===================================================================
\section{Classification-Surface, Cotenancy, and Extraction Validation}
\label{sec:empirical}

\subsection{Setup}

The formative classification experiment applies a regime mask to the final
pooled encoder vector before a shared classifier; it does not place
the mask inside each Transformer FFN.  We report them as formative
evidence that gate-aware optimization can learn useful
support-constrained representations, not as empirical FFN-placement
evidence.  DistilBERT is evaluated on AG
News~\citep{zhang2015agnews} four-class classification
(8,000--20,000 training examples and 7,600 test examples).  We then
separately test standalone extraction of gated projection columns.
Paired statistical tests use five seeds where explicitly reported.
The strict cotenancy experiments below instead gate every expanded
intermediate FFN and use three seeds unless stated otherwise.

\begin{table}[ht]
\centering
\caption{Post-encoder classification configurations.}
\label{tab:models}
\begin{tabular}{@{}llrrl@{}}
\toprule
\textbf{Model} & \textbf{Hidden dim} & \textbf{Params} &
\textbf{R tested} & \textbf{Training} \\
\midrule
DistilBERT-base & 768 & 66M & 8--128 & Co-trained \\
\bottomrule
\end{tabular}
\end{table}

\subsection{Co-Training Cost Versus Separate Controls}
\label{sec:series1}

The updated DistilBERT sweep uses five paired seeds at each of seven
partition ratios on parallel A100 workers.  In this formative
protocol, the mask acts on the 768-dimensional final CLS vector;
``Dims/Regime'' below therefore refers to pooled-classifier input
coordinates, not DistilBERT's 3,072-dimensional intermediate FFN.
At every ratio, the gated model is trained across all $R$ regimes,
whereas the separate-model control samples four regimes; each
per-seed statistic pairs the mean of those four controls with the
corresponding four gated-regime accuracies.

\begin{table}[ht]
\centering
\caption{Five-seed paired co-training results (DistilBERT, AG News).
Gap is separate minus gated accuracy in percentage points.  Intervals
use Student's $t$ distribution with four degrees of freedom.}
\label{tab:series1}
\small
\begin{tabular}{@{}rrrrrr@{}}
\toprule
$R$ & Dims/Regime & Gap (pp) & 95\% CI (pp) & $p$ & Bonf.\ reject? \\
\midrule
8   & 96 & 0.348 & [$-0.229$, 0.925] & 0.1692 & No \\
16  & 48 & 0.268 & [$-0.301$, 0.837] & 0.2616 & No \\
24  & 32 & 0.208 & [$-0.324$, 0.740] & 0.3392 & No \\
32  & 24 & 0.267 & [$-0.345$, 0.879] & 0.2920 & No \\
64  & 12 & 0.143 & [$-0.041$, 0.326] & 0.0967 & No \\
96  & 8  & 0.376 & [0.045, 0.708]    & 0.0345 & No \\
128 & 6  & 0.272 & [$-0.208$, 0.752] & 0.1902 & No \\
\bottomrule
\end{tabular}
\end{table}

Across the seven ratios, the measured cost is small: the
separate-minus-gated gap ranges from 0.143 to 0.376 percentage
points.  The $R{=}96$ paired test is significant at the unadjusted
0.05 level, but no test crosses the Bonferroni family-wise threshold
$0.05/7 \approx 0.0071$.  Bonferroni is intentionally conservative;
we include it because seven related hypotheses were inspected.  The
ratios share the same paired seed schedule and are correlated, and
the $R{=}96$ result is treated as an exploratory post-hoc observation
rather than a preregistered discovery.
Failure to reject a zero difference is not proof of equality, nor is
exact equality expected from independently optimized inscriptions.
The defensible conclusion is that the observed performance costs are
small in this five-seed DistilBERT study.  A formal equivalence claim
would additionally require a predeclared practical margin and an
interval wholly contained within that margin.

The archived experiment initially printed normal-approximation
intervals using $1.96$ despite only five pairs.  \Cref{tab:series1}
recomputes the intervals with the matching Student-$t$ critical value;
this correction also removes an apparent inconsistency between the
$R{=}64$ interval and its paired-test $p$-value.

\begin{sloppypar}
\noindent\textit{Replay.} Use the legacy post-encoder surface in
\path{experiments/reproduce_distilbert_classification.py}; the reported
five-seed rows are retained in
\path{experiments/results/series1_results_20260803_202919.json}.
\end{sloppypar}

\subsection{Holding Per-Regime Classifier Capacity Fixed}
\label{sec:capacity-preserving-classifier}

The preceding sweep fixes the 768-dimensional post-encoder
representation, so each support narrows as $R$ grows.  This is not the
only sizing policy.  If a task needs $q$ hidden coordinates per regime,
the governed layer may instead choose $d=qR$.  Adding a regime then
adds a block without reducing any incumbent regime's allocated width.

We ported and reran the formative wide-classifier protocol with
$R{=}128$, $q{=}10$, and therefore $d{=}1{,}280$.  The two-layer MLP
has 256 one-hot synthetic observations and 256 disjoint answers, two
per regime.  It is co-trained for 2,000 full-batch epochs using a
post-activation binary gate.  The sparse selected-support
implementation is also compared with explicit dense-mask execution.

The owning key produces the expected answer for 256/256 observations.
Across all 32,512 pairs of a foreign observation and wrong regime key,
the expected foreign answer is never produced.  Activating all 128
regimes simultaneously produces 249/256 correct answers, and
selected-support versus dense-mask logits have maximum absolute
difference zero.

This result is deliberately narrow.  It is a deterministic synthetic
memorization and scalability test with no held-out set, and the zero
wrong-key hit count is an empirical negative control rather than a
privacy theorem.  It does show why ``$R$ regimes'' need not imply
``$1/R$ of a fixed model'' for modular classifiers: each regime retains
the same ten-coordinate allocation because the governed layer grows
with $R$.  The cost is linear parameter storage, not lost per-regime
width.

\begin{sloppypar}
\noindent\textit{Replay.} The self-contained runner is
\path{experiments/capacity_preserving_wide_classifier.py}; the reported
run is
\path{experiments/results/capacity_preserving_wide_20260820_225331.json}.
\end{sloppypar}

\subsection{Deployment Service-Consolidation Consequence}
\label{sec:service-consolidation}

An $R{=}8$ DistilBERT SST-2 deployment benchmark compares eight complete model
services with three single-backbone layouts.  Eight services require 1.071 GB
of checkpoint tensors and 1.103 GB of resident CUDA parameter and buffer state.
A frozen shared backbone with eight zero-initialized private adapter slots
retains the original 91.06\% accuracy using 151.6 MB and 156.2 MB,
respectively.  This is a deployment control, not evidence that untrained
adapters add utility.

Post-hoc one-eighth FFN slicing reduces accuracy to 67.66\%.  Physically
extracting the authorized slices preserves the sliced logits within the
declared fp16 tolerance and raises measured throughput from 1,585 to 2,872
samples/s.  Every scored request passes through execution authority and malformed
authority makes zero model calls.  Thus the measurement supports consolidating
a frozen backbone plus private modules, while directly rejecting the stronger
claim that naive narrowing of a trained dense Transformer preserves utility.
Timing is one serial stream on one NVIDIA T4 and is not a concurrency result.

\begin{sloppypar}
\noindent\textit{Replay.} Run
\path{experiments/modal_distilbert_service_consolidation.py}; the full artifact
is listed below.
\par\smallskip\noindent{\footnotesize
\path{experiments/results/distilbert_service_consolidation_20260821T135530_707107Z.json}.}
\par
\end{sloppypar}

\subsection{Strict Frozen-Backbone Transformer Cotenancy}
\label{sec:strict-cotenancy}

The formative post-encoder study above establishes task feasibility,
not a separated shared-model state boundary: its optimizer may update
attention, normalization, and other shared encoder parameters.  We
therefore run a second protocol whose acceptance condition is stricter
than task accuracy.  DistilBERT's attention, embeddings,
normalization, residual parameters, and all other non-governed shared
weights are frozen.  A gate is placed on the 3,072-dimensional
expanded activation immediately before every FFN down projection.
For regime $r$, training may update only the active rows of
$\bW_1$, the matching coordinates of $\mathbf{b}_1$, the matching
columns of $\bW_2$, and a private regime-selected classifier in
PyTorch's \texttt{Linear.weight} storage convention.  Those tensors
are transposes of the row-vector notation in
\eqref{eq:intermediate-ffn-gate}, where the same slices are columns of
$\bW_1$ and rows of $\bW_2$.
A regime-scoped Adam implementation updates moments only on those
authorized entries; ordinary Adam momentum would otherwise continue
moving a previously active coordinate during a later tenant's step
even when its current gradient is zero.

We train on 8,000 AG News examples for two epochs and evaluate on
2,000 held-out examples at seeds 42, 123, and 256.  Label
permutations provide regime-specific functional identities.  The
$R{=}1$ condition is the full-FFN-capacity reference under the same
frozen-backbone protocol.

\begin{table}[ht]
\centering
\caption{Strict intermediate-FFN cotenancy on frozen DistilBERT.
Accuracy is mean $\pm$ sample standard deviation across three seeds.
The drop is relative to the matched $R{=}1$ capacity reference.}
\label{tab:strict-dense-ffn}
\small
\begin{tabular}{@{}rrrrr@{}}
\toprule
$R$ & FFN dims/tenant/layer & Owning accuracy & Drop (pp) &
Max unauthorized delta \\
\midrule
1  & 3,072 & $91.28\%\pm0.33\%$ & 0.00 & 0 \\
2  & 1,536 & $90.15\%\pm0.39\%$ & 1.13 & 0 \\
4  &   768 & $88.93\%\pm0.30\%$ & 2.36 & 0 \\
8  &   384 & $87.47\%\pm0.04\%$ & 3.82 & 0 \\
16 &   192 & $86.21\%\pm0.08\%$ & 5.07 & 0 \\
\bottomrule
\end{tabular}
\end{table}

``Max unauthorized delta'' is the maximum over frozen shared
parameters, off-partition FFN parameters, off-partition first and
second optimizer moments, and inactive private classifiers in all
runs.  Every category is bit-exact zero.  Thus the capacity cost is
empirical and monotone in this setup, while the declared state
boundary is exact.  Halving each FFN retains 90.15\% accuracy,
1.13 percentage points
below the full-capacity reference.  At $R{=}16$, the 192-coordinate
slice incurs a material but non-catastrophic 5.07-point drop.  This is
a task- and architecture-specific capacity curve, not a universal
Transformer scaling law.

\begin{sloppypar}
\noindent\textit{Replay.} Run
\path{experiments/modal_dense_ffn_cotenancy.py}; the three-seed
artifacts are the
\path{experiments/results/dense_ffn_cotenancy_*.json} files and their
aggregate is
\path{experiments/results/transformer_cotenancy_summary.json}.
\end{sloppypar}

\paragraph{Preliminary public gate-aware adaptation.}
The preceding sweep applies intermediate-FFN masks post hoc to the
pretrained backbone.  We separately test whether a backbone can learn
the $R{=}8$ partition geometry without receiving tenant data.  At each
of seeds 42, 123, and 256, an ungated teacher is fine-tuned for two
epochs on 8,000 examples from a designated public AG News split.  A
student copied from that teacher is then trained for one epoch on the
same public split; for every batch, it is evaluated through each of
the eight 384-coordinate FFN masks and optimized with an equal mixture
of hard-label loss and temperature-2 teacher distillation loss.
The comparison checkpoint retains the teacher without this additional
epoch and is therefore a no-extra-training baseline, not an
equal-compute control.  The next 8,000 shuffled training examples form a disjoint
tenant-stage split.  For both conditions, all shared paths are then
frozen and only aligned FFN slices and private classifiers are trained
for two epochs, exactly as in the strict protocol above.  The first
2,000 standard test examples are held out from both stages.

\begin{table}[ht]
\centering
\caption{Preliminary $R{=}8$ public-adaptation comparison.  Accuracy is
mean $\pm$ sample standard deviation across three seeds.  Uplift is
relative to the no-extra-training baseline.}
\label{tab:public-gate-adaptation}
\small
\begin{tabular}{@{}lrrr@{}}
\toprule
Initialization & Owning accuracy & Difference (pp) & Max delta \\
\midrule
Public task only & $88.55\%\pm0.11\%$ & 0.00 & 0 \\
$+$ all-mask distillation & $90.15\%\pm0.11\%$ &
$1.60\pm0.10$ & 0 \\
\bottomrule
\end{tabular}
\end{table}

Before tenant training, the ungated teacher reaches
$91.62\%\pm0.35\%$.  Applying the eight masks to the generic-public
control yields $74.60\%\pm5.42\%$ mean accuracy, while the all-mask
student reaches $89.89\%\pm0.51\%$.  After the strict tenant stage,
the combined extra-training, masking, and distillation pipeline retains
the 1.60-point difference and
finishes 1.47 points below the ungated teacher.  Frozen shared
parameters, off-partition FFN parameters and optimizer moments, and
inactive classifiers all remain bit-exact zero during tenant
training.  Because the baseline did not receive the extra public
epoch, this comparison cannot identify whether the difference comes
from additional training, mask-aware adaptation, distillation, or
their interaction.  It is retained as preliminary pipeline evidence,
not a causal estimate and not evidence of zero utility loss.  The
``public'' and ``tenant'' designations are a controlled disjoint split
of one public benchmark, and label permutations supply functional
identities; real private-corpus generalization remains outside this
test.

\begin{sloppypar}
\noindent\textit{Replay preliminary adaptation.} The reported
three-seed artifact is
\path{experiments/results/public_gate_distillation_20260806T193105_795325Z.json}.
The repository retains the stronger equal-work successor below, but
not the exact historical runner that produced this artifact.
\end{sloppypar}

\paragraph{Equal-work factorial result.}
The checked-in successor runner copies the same public teacher into
four conditions: no extra public adaptation, an extra ungated
hard-label epoch, an extra all-mask hard-label epoch, and an extra
all-mask hard-label-plus-distillation epoch.  The three adapted
conditions use the same examples, batch order, optimizer steps,
teacher evaluations, and student forward/backward counts; the
ungated arm repeats each batch $R$ times to match model-pass counts.
The resulting paired contrasts estimate the effects of extra
training, mask-aware training, and distillation conditional on masks.
In the completed full seed, ordinary extra training changes tenant accuracy by
$+0.075$ percentage points, mask awareness adds $+1.331$ points beyond that
control, and distillation adds $+0.031$ points conditional on masks. The full
pipeline is $+1.438$ points above no extra adaptation, with exact zero frozen,
off-partition, optimizer-moment, and inactive-classifier deltas. The reduced
pilot is negative and retained as a protocol check. One full seed does not
support a population-level effect claim.

\begin{sloppypar}
\noindent\textit{Replay factorial.} Run
\path{experiments/modal_public_gate_adaptation_factorial.py}; the completed
artifact is
\path{experiments/results/public_gate_adaptation_factorial_20260821T063703_914999Z.json}.
\end{sloppypar}

\paragraph{Complete private lanes on shared attention.}
We also test two designs that avoid dividing a trainable tenant module
by $R$.  Both keep the pretrained DistilBERT backbone frozen and
shared.  The \emph{adapter} design adds a 64-dimensional private
residual adapter after every Transformer block; the \emph{expert}
design adds a full 3,072-dimensional private FFN after every block.
The regime-selected classifier is private in both designs.  The
private output then enters the next frozen shared attention block:
attention is useful and unablated, but remains inside the trusted
shared boundary.

\begin{table}[ht]
\centering
\caption{Complete private lanes on a frozen shared DistilBERT
backbone ($R{=}4$).  Accuracy is mean $\pm$ sample standard deviation
across three seeds.}
\label{tab:private-lanes}
\small
\begin{tabular}{@{}lrrr@{}}
\toprule
Design & Owning accuracy & Wrong-key accuracy &
Max shared/inactive delta \\
\midrule
64-dim residual adapters & $90.47\%\pm0.18\%$ &
$3.18\%\pm0.06\%$ & 0 \\
3,072-dim private experts & $90.09\%\pm0.14\%$ &
$3.30\%\pm0.05\%$ & 0 \\
\bottomrule
\end{tabular}
\end{table}

The private adapter is the best measured tradeoff here: it slightly
outperforms the much larger expert while adding approximately
0.59M trainable parameters per tenant rather than 28.33M.  The
wrong-key values are functional negative controls induced by distinct
label permutations; they are not the source of the security claim.
The exact frozen-backbone and inactive-lane deltas establish that
tenant training state does not enter another lane under the declared
optimizer and runtime boundary.  Shared activations are not claimed
to be tenant-private, and cache, batching, and serving-runtime state
must be separately namespaced.

\begin{sloppypar}
\noindent\textit{Replay.} Run
\path{experiments/modal_private_transformer_lanes.py}; the reported
adapter and expert runs are
\path{experiments/results/private_transformer_lanes_*.json}.
\end{sloppypar}

\paragraph{Key-authorized customer corpora and Cargo Mode.}
\label{sec:cargo-transformer}
Finally, a legacy two-tenant smoke test connects the model-state
boundary to runtime data authorization.  Each customer corpus occupies
a separate vector-store partition.  The lockbox side retains the
authorization root key and grants a scoped Cargo access key.  The
completed artifact binds tenant and regime but predates the canonical
sub-id requirement.  It is retained as a real-model mechanism smoke,
not as evidence that the current complete scope was enforced.  Its
manifest also binds the embedding
specification, payload hash, and item count; payload count and hash
are validated before ingest.  Retrieval must successfully dock at the
corresponding regime bus before context is passed to an otherwise
normal shared FLAN-T5 generator.

All four owning questions retrieve the correct customer partition and
produce the exact canary answer.  Six unauthorized attempts
(random key, other-tenant key, and other-regime key against each
tenant) are rejected before retrieval, with zero unauthorized model
calls.  All load and retrieval receipts verify under the scoped access
keys.  For this retained legacy artifact, ``verify'' denotes the recorded
signature check; it is not evidence that the receipt committed to the exact
returned content.  These small counts are protocol checks, not statistical
leakage estimates.  The enforcement claim depends on the vector store, master
key, lockbox proxy, and runtime process remaining trusted and on
corpus access being exposed only through the regime bus.  This design
lets ordinary Transformer attention consume authorized customer
context; it does not partition or ablate attention heads.

The checked-in successor protocol is self-contained and enforces the
complete research scope.  It binds tenant, sub-id, regime, model digest, operation,
policy version, and the registered vector partition, with distinct load
and retrieval keys plus wrong-sub-id, wrong-operation, and same-regime
wrong-partition negative controls.  It now also places a separate
tenant + sub-id + model + operation gate before generator invocation.
It is configured as a protocol rerun; no result from that strengthened
GPU runner is claimed until its artifact is checked in.  The completed
local CPU protocol check exercises the same fail-closed scope logic,
including wrong-sub-id and whole-model negatives with zero unauthorized
model calls, without making a generative-quality claim.

\begin{sloppypar}
\noindent\textit{Replay.} Run
\path{experiments/local_transformer_cotenancy_suite.py} for the Gate-backed local
scope matrix or
\path{experiments/modal_cargo_transformer_authorization.py} for the
real-model Cargo experiment.  Retained outputs and their evidentiary
status are indexed by \path{experiments/results/README.md}.
\end{sloppypar}

\subsection{Structural FFN Cross-Regime Coordinate Exclusion}
\label{sec:isolation}

For $r \neq s$, disjointness gives
$\bg_r \odot \bg_s=\mathbf{0}$: a runtime invocation resolved to
regime $s$ cannot activate coordinates owned by regime $r$.  The
lockbox generates the global Fisher--Yates assignment and scoped
capabilities; the trusted runtime performs credential-to-mask
dispatch.  Callers do
not choose a permutation directly.  This coordinate exclusion is
the structural guarantee.  The formative runner also computes
wrong-key task accuracy, but the checked-in parity-only artifact has
an empty scaling block and therefore does not contain the per-regime
or cross-regime measurements previously quoted here.  We make no
numerical wrong-key claim from that formative artifact.  The distinct
strict-cotenancy runs in \Cref{sec:strict-cotenancy} do archive
wrong-key functional controls, but their security statement rests on
exact inactive state and runtime rejection rather than low task
accuracy.  The exact coordinate zero does not depend on either
task-level measurement.

\subsection{Deployment Efficiency}
\label{sec:deployment}

A shared gated deployment stores one backbone plus masks and any
private heads or lanes, whereas a baseline that materializes $R$
identical architecture copies stores the backbone $R$ times.  At
$R{=}4$, the checked-in CPU smoke reproduces the resulting $4\times$
raw state-dict arithmetic for the identical-copy baseline; it does
not establish utility equivalence or complete production memory
savings.  A public runner
also measures simultaneous accelerator allocation and routed versus
multiplexed throughput, but the repository does not yet contain a
full GPU artifact supporting the historical $3.6\times$ memory and
$2.69\times$ throughput figures.  Those figures are withheld pending
a hardware-matched rerun.  The benchmark applies per-sample masks at
the pooled classifier input and is not an intermediate-FFN systems
benchmark.

\begin{sloppypar}
\noindent\textit{Replay measurements.} Storage, live accelerator
allocation, and routed-throughput collection are implemented in
\path{experiments/reproduce_distilbert_deployment.py}; the retained
publication-safe artifact is
\path{experiments/results/distilbert_deployment_20260806_103232.json}.
\end{sloppypar}

\subsection{Exact Submodel Extraction}
\label{sec:extraction}

This experiment asks whether an authorized regime can be exported as
a smaller standalone linear projection without changing the function
computed by the full gated projection.  For a full matrix
$\bW\in\mathbb{R}^{4\times768}$, input batch
$\mathbf{X}\in\mathbb{R}^{256\times768}$, bias $\mathbf{b}$, and active
coordinate set $A_r$, we compare
\begin{equation}
  (\mathbf{X}\odot\bg_r)\bW^\top+\mathbf{b}
  \quad\text{with}\quad
  \mathbf{X}[:,A_r]\,\bW[:,A_r]^\top+\mathbf{b}.
\end{equation}
The left side executes the full matrix after applying the gate; the
right side physically removes every inactive input coordinate and
executes only the extracted columns.  Equality is expected
algebraically because all omitted products are multiplication by
zero.  The empirical test determines whether implementation details
or floating-point reduction order create a numerical difference.

\paragraph{Strategy.}
For each of 10 random seeds, we generate a random input batch, weight
matrix, bias, and permutation of 768 coordinates.  The permutation is
split into four disjoint 192-coordinate regimes.  We test each regime
individually and then test authorized unions containing one, two,
three, and all four regimes.  The complete matrix is repeated in
32-bit and 16-bit floating point, yielding 80 single-regime checks and
80 compound-mask checks.  The predeclared absolute tolerances are
$10^{-6}$ for 32-bit and $2\times10^{-2}$ for 16-bit, with an
additional 16-bit relative tolerance of $10^{-3}$.

\begin{table}[ht]
\centering
\caption{Exact-extraction matrix.  ``Bit-equal'' counts comparisons
whose output tensors have identical bits; all 160 comparisons satisfy
the declared numeric tolerance.}
\label{tab:extraction-matrix}
\begin{tabular}{@{}lrrrr@{}}
\toprule
Precision & Comparisons & Bit-equal & Within tolerance & Max $|\Delta|$ \\
\midrule
32-bit floating point & 80 & 80 & 80 & 0 \\
16-bit floating point & 80 & 60 & 80 & 0.015625 \\
\bottomrule
\end{tabular}
\end{table}

The result confirms exact execution in 32-bit arithmetic for this
matrix and tolerance-bounded equivalence in 16-bit arithmetic.  It
matches the Lean~4 gated-compressed equivalence theorem
(\texttt{compressed\_eq\_gated}), while making the finite-precision
qualification explicit.  For one of $R$ equal-width regimes, the
extracted projection retains $1/R$ of the gated input-facing weight
columns; authorized compound masks retain the union of their columns.
This is also the linear algebra of extracting the active input rows of
$\bW_1$ in the row-vector pre-MLP convention.  The separate local FFN
protocol check instantiates a $32\mathbin{\to}64$ first projection and
observes zero inactive input and $\bW_1$-row gradients.  Neither check
is a pre-MLP task-utility benchmark.

\begin{sloppypar}
\noindent\textit{Replay.} Run
\path{experiments/local_exact_extraction.py}; the reported artifact is
\path{experiments/results/local_exact_extraction_20260803_172654.json}.
\end{sloppypar}

\subsection{Excluded Experimental Scope}
\label{sec:excluded-experiments}

The empirical claims in this section include post-encoder
classification, strict intermediate-FFN and private-lane cotenancy,
the keyed Cargo retrieval smoke, checked-in storage, exact extraction,
and local pre-MLP/whole-model authorization protocol checks.
Exploratory post-pooled embedding, open-ended generative
tenant-state training, complete attention-lane partitioning,
adapter-state exfiltration, and KV-cache studies remain excluded:
they either do not instantiate the declared intermediate FFN/runtime
boundary or address systems beyond this paper's scope.  Their
artifacts remain archived for future work but are not used as
evidence here.

% ===================================================================
\section{Supplementary Keyed Permutation Conjugation}
\label{sec:orthogonal-superposition}

The previous sections partition a model's hidden dimensions into
disjoint subspaces using binary gate masks.  A complementary
question arises: can a single trained function be assigned $R$
keyed basis placements
by exploiting the fact that a transformer's computation is
equivariant under permutation of the residual-stream basis?  The
permutation creates the keyed placement, not tenant knowledge or an
additional isolation boundary.

\subsection{Permutation Equivariance of Transformers}
\label{sec:permutation-equivariance}

\begin{theorem}[Residual-stream permutation equivariance]
\label{thm:permutation-equivariance}
Let $\mathcal{T}$ be a transformer with residual-stream
dimension~$d$, and let $P \in \{0,1\}^{d \times d}$ be a
permutation matrix.  Write residual activations as row vectors and
relabel them by $\bx\mapsto\bx P^\top$.  Define the \emph{conjugated} transformer
$\mathcal{T}^{(P)}$ by replacing every residual-facing parameter:
\begin{align}
    \bW_{\mathrm{emb}} &\mapsto \bW_{\mathrm{emb}} P^\top
    & \text{(embedding outputs)} \label{eq:conj-emb} \\
    \bW_Q, \bW_K, \bW_V &\mapsto P\bW_Q, \;P\bW_K,\;
    P\bW_V
    & \text{(attention inputs)} \label{eq:conj-qkv} \\
    \bW_O &\mapsto \bW_O P^\top
    & \text{(attention output)} \label{eq:conj-out} \\
    \bW_1^{\mathrm{FFN}} &\mapsto P\bW_1^{\mathrm{FFN}}
    & \text{(FFN input)} \label{eq:conj-ffn1} \\
    \bW_2^{\mathrm{FFN}} &\mapsto \bW_2^{\mathrm{FFN}}P^\top
    & \text{(FFN output)} \label{eq:conj-ffn2} \\
    \boldsymbol{\gamma}, \boldsymbol{\beta} &\mapsto
    \boldsymbol{\gamma}P^\top, \; \boldsymbol{\beta}P^\top
    & \text{(LayerNorm)} \label{eq:conj-ln} \\
    \mathbf b_O,\mathbf b_2 &\mapsto
    \mathbf b_OP^\top,\;\mathbf b_2P^\top
    & \text{(residual-output biases)} \label{eq:conj-bias} \\
    \bW_{\mathrm{cls}} &\mapsto P\bW_{\mathrm{cls}}
    & \text{(classifier input)} \label{eq:conj-cls}
\end{align}
Then $\mathcal{T}^{(P)}$ produces identical outputs to
$\mathcal{T}$ on all inputs:
\begin{equation}
    \mathcal{T}^{(P)}(\bx) = \mathcal{T}(\bx)
    \quad \forall\, \bx
    \label{eq:equivariance}
\end{equation}
\end{theorem}

\begin{proof}[Proof sketch]
The conjugation relabels the residual-stream coordinate system.
At every input-side map, $\bx P^\top(P\bW)=\bx\bW$; every
residual-output map appends $P^\top$.  LayerNorm's mean and variance
are permutation-invariant, and its scale and shift are relabeled by
the same $P^\top$, so LayerNorm commutes with the conjugation.
Attention computes $\bQ\bK^\top$ on the \emph{per-head}
subspace, which is an internal coordinate system unaffected by
the residual-stream permutation once $\bW_Q$, $\bW_K$, $\bW_V$
absorb $P$ on their input side.  The result is a change
of basis that is invisible to the model's computation.
\end{proof}

The theorem implies that for any trained model and any $R \leq d!$
distinct permutations $\{P_0, P_1, \ldots, P_{R-1}\}$, the $R$
conjugated models $\{\mathcal{T}^{(P_r)}\}$ are functionally
identical yet operate in mutually distinct coordinate systems.
Each conjugated model is a keyed \emph{basis placement}: the
accuracy gap from the base model is exactly zero, not approximately
zero.  Conjugation alone does not create distinct knowledge or
zero non-owning coordinates.  A disjoint CDP mask can provide the
coordinate zeros in the selected basis; tenant-specific deposition
into those coordinates is not evaluated here.

\paragraph{Relationship to CDP.}
Permutation conjugation and CDP's Hadamard gate are complementary
mechanisms.  CDP partitions dimensions via binary masks to
provide FFN-coordinate confinement.  Permutation
conjugation relabels the entire coordinate system to achieve
\emph{keyed basis placement} (different regimes use different
coordinate systems for the same initial computation).  The two compose
naturally for a single selected basis: an outer permutation
conjugation places the model in that orientation, and an inner CDP
gate selects FFN coordinates within it.  This composition preserves
the gate theorem but does not itself establish tenant-specific
learning.
The credential, custody, and dispatch assumptions required for this
composition are stated in \Cref{sec:threat-model}.

Operationally, the inner gate and outer shuffle perform
different jobs.  The binary mask first determines which FFN
coordinates are active;
the residual-stream permutation then removes a stable public
labeling of the surrounding representation and binds its
realization to the engine-held key.  For an FFN gate, $P_r$ is
absorbed at the residual-facing sides of $\bW_1$ and $\bW_2$,
while $\bg_r$ acts on the internal FFN activation between them.
Thus the two operations do not compete: ablation determines
the active FFN support, and the keyed shuffle determines
the coordinate system in which the trusted runtime realizes it.
If a gate and permutation act on the same coordinate space, the
equivalent regime-local mask is simply $P_r\bg_r$.

The reference code contains the conjugation, multiplexing,
paper-aligned intermediate FFN gate, and nested-mask paths.  The
algebraic composition is supported by
\texttt{permute\_hadamard} and
\texttt{permutation\_preserves\_logits} in Lean~4.  These paths are
not yet packaged as one unified production execution mode; the
security claim for a deployment follows the specific composed path
it actually instantiates.

\subsection{Empirical Validation}
\label{sec:superposition-empirical}

We validate residual-stream permutation conjugation on
DistilBERT-base~\citep{sanh2019distilbert} (66M parameters,
$d{=}768$, 6 layers, 12 heads) fine-tuned on AG News~\citep{zhang2015agnews}
4-class classification (8,000 train / 7,600 test, 4~epochs,
AdamW $\mathrm{lr}{=}2{\times}10^{-5}$, seed~0).  All
experiments run on Apple MPS (M-series laptop).  Base accuracy:
90.25\%.

\paragraph{Experiment A: R-sweep.}
Train one model.  For each $R \in \{4, 8, 16, 32, 64, 128\}$,
generate $R$ random permutations (seeded), deep-copy the model,
conjugate all residual-facing parameters per
\Cref{thm:permutation-equivariance}, and evaluate on the full
test set.  At $R{=}1024$, sample 8 regimes across the full
range ($r \in \{0, 1, 10, 100, 500, 777, 999, 1023\}$).

\begin{table}[ht]
\centering
\caption{Keyed permutation-conjugation R-sweep.  Accuracy gap is
exactly zero at every $R$ tested, from 4 to 1,024.}
\label{tab:r-sweep}
\small
\begin{tabular}{@{}rrrr@{}}
\toprule
$R$ & Regimes tested & Max $|\Delta\mathrm{acc}|$
& Accuracy \\
\midrule
4 & 4 & $0.00$ & 90.25\% \\
8 & 8 & $0.00$ & 90.25\% \\
16 & 16 & $0.00$ & 90.25\% \\
32 & 32 & $0.00$ & 90.25\% \\
64 & 64 & $0.00$ & 90.25\% \\
128 & 128 & $0.00$ & 90.25\% \\
1{,}024 & 8 (sampled) & $0.00$ & 90.25\% \\
\bottomrule
\end{tabular}
\end{table}

The measured accuracy gap is zero at every regime tested, including
regime indices 777 and 1023 at $R{=}1024$.  This prediction-level
result is consistent with the theorem; exact functional identity
follows from the conjugation algebra rather than from finite test-set
accuracy alone.

The experiment tests an expected invariance of a single function; it
does not establish additional learned functions, tenant state, or an
isolation boundary.  A permutation vector is only addressing metadata
for that reparameterization, so comparison with materialized copies
of the same function is not presented as model compression.

\paragraph{Experiment B: component ablation.}
To verify that conjugation must be total, we systematically
skip conjugation for each component class (embeddings, LayerNorm,
attention projections, FFN weights, classifier) and measure
the resulting accuracy change.

\begin{table}[ht]
\centering
\caption{Component ablation: skipping each component from
conjugation changes accuracy, often severely.  Base accuracy 90.25\%.}
\label{tab:ablation}
\small
\begin{tabular}{@{}lrr@{}}
\toprule
\textbf{Skipped component} & \textbf{Accuracy} & \textbf{Gap} \\
\midrule
None (full conjugation) & 90.25\% & $0.00$ \\
\midrule
Embeddings & 24.72\% & $-65.53$ \\
Embedding LayerNorm & 88.59\% & $-1.66$ \\
SA LayerNorm & 39.89\% & $-50.36$ \\
Output LayerNorm & 43.41\% & $-46.84$ \\
All LayerNorm & 41.11\% & $-49.14$ \\
Q projection & 81.95\% & $-8.30$ \\
K projection & 77.07\% & $-13.18$ \\
V projection & 21.86\% & $-68.39$ \\
All Q/K/V & 21.68\% & $-68.57$ \\
Out projection ($\bW_O$) & 24.96\% & $-65.29$ \\
FFN $\bW_1$ & 25.00\% & $-65.25$ \\
FFN $\bW_2$ & 52.21\% & $-38.04$ \\
FFN (both) & 27.80\% & $-62.45$ \\
Classifier & 45.78\% & $-44.47$ \\
\bottomrule
\end{tabular}
\end{table}

In this implementation, omitting every tested component changes
accuracy.  The most destructive single omission in this run is the
V-projection (collapse to 21.86\%, below the 25\% chance baseline),
but the ablation does not establish a general ranking of information
paths.  Even the embedding LayerNorm causes a 1.66\,pp gap when
skipped, because its element-wise scale and shift parameters are
residual-facing.

\subsection{Keyed-Basis Batched Inference}
\label{sec:true-multiplexing}

Permutation conjugation can batch $R$ keyed basis placements through
one set of weights.  This is a systems implementation of
function-preserving reparameterization, not tenant isolation or a new
batching primitive.  Instead of
pre-conjugating model copies, we apply per-sample permutations to
the activations at every residual-stream boundary via
\texttt{torch.gather}:

\begin{enumerate}[nosep]
    \item Before each weight matrix (input side): inverse-permute
    the activations to the original basis so the shared weights
    apply correctly.
    \item After each weight matrix (output side): forward-permute
    the result back to the sample's regime basis.
    \item Around each LayerNorm: inverse-permute, normalize, then
    forward-permute (because LayerNorm parameters are
    element-wise in the original basis).
\end{enumerate}

The weight matrices are never modified; only the activations are
shuffled per-sample.  Correctness was verified at $R{=}4$: the
maximum logit difference between keyed-basis batched and serial
(pre-conjugated copy) evaluation is $1.19 \times 10^{-6}$, with
bit-exact predictions.

\begin{table}[ht]
\centering
\caption{Keyed-basis batching: time to evaluate $R$ samples
(one per basis placement) through a single DistilBERT backbone.
Apple MPS, 30 trials, 5 warmup.}
\label{tab:mux-scaling}
\begin{tabular}{@{}rrrrrr@{}}
\toprule
$R$ & Keyed batch & Serial & Speedup
& Per-sample & Cost vs.\ $R{=}1$ \\
\midrule
1 & 2.79\,ms & 2.66\,ms & $0.95\times$ & 2.79\,ms & $1.00\times$ \\
4 & 6.28\,ms & 10.73\,ms & $1.71\times$ & 1.57\,ms & $2.25\times$ \\
8 & 10.92\,ms & 39.24\,ms & $3.59\times$ & 1.37\,ms & $3.91\times$ \\
16 & 21.92\,ms & 45.53\,ms & $2.08\times$ & 1.37\,ms & $7.86\times$ \\
32 & 43.35\,ms & 94.60\,ms & $2.18\times$ & 1.35\,ms & $15.54\times$ \\
64 & 90.66\,ms & 205.44\,ms & $2.27\times$ & 1.42\,ms & $32.49\times$ \\
128 & 339.03\,ms & 469.94\,ms & $1.39\times$ & 2.65\,ms & $121.52\times$ \\
\bottomrule
\end{tabular}
\end{table}

The per-sample cost drops from 2.79\,ms ($R{=}1$, no permutation)
to a floor of 1.35\,ms ($R{=}32$) as GPU utilization improves with
larger batch sizes.  At $R{=}128$, the per-sample cost rises to
2.65\,ms as memory bandwidth saturates on the test hardware.  The
total cost of evaluating 128 keyed basis placements ($339$\,ms) is
$121.5\times$ the cost of one.  Ordinary batching improves
accelerator utilization while the permutation gathers add per-sample
work.  The benchmark does not include an unpermuted batch at the same
sizes, so it does not isolate permutation overhead from ordinary
batching efficiency.  It is a correctness and placement benchmark,
not evidence of tenant-specific deposition, isolation, or a new
batching primitive.

\paragraph{Idealized operation-count ratio.}
The gather operations contribute $\mathcal{O}(B \cdot
\mathrm{seq} \cdot d)$ memory bandwidth per layer.  The matrix
multiplications contribute $\mathcal{O}(B \cdot \mathrm{seq}
\cdot d^2)$.  The ratio is $\mathcal{O}(1/d) \approx 1/768
\approx 0.13\%$.  This ratio omits kernel launch, memory movement,
indexing, and hardware effects and is not a measured latency floor.
The measured overhead (93.8\% at $R{=}128$
on MPS) is dominated by Python-level \texttt{torch.gather}
calls at each injection point.  A hypothetical fused kernel that
permutes weight-matrix indices during the GEMM---or
precomputes permuted weight views ($\bW[:,P]$)---could reduce this
overhead, but no such kernel is implemented or benchmarked here.

\begin{sloppypar}
\noindent\textit{Replay.} Run
\path{experiments/orthogonal_superposition_sweep.py} for permutation
conjugation and \path{experiments/true_multiplexing_test.py} for keyed
basis batching.  Retained sweep, ablation, and timing outputs are
indexed by \path{experiments/results/README.md}.
\end{sloppypar}

% ===================================================================
\section{Discussion}
\label{sec:discussion}

\paragraph{CDP as measurement instrument.}
End-to-end cross-tenant leakage is ordinarily estimated
statistically---for example with distribution comparisons,
membership-inference attacks, or accuracy deltas.  CDP adds a local
structural zero at its declared FFN gate.
The provable baseline is local and exact: an inactive coordinate at
the declared gate is zero, so no signal can traverse that coordinate
($0 \cdot x = 0$).  The reported 0.499~AUC membership-inference
score and 0/30 factual exfiltration result belong to prior work and
are not evidence in this paper.  Previously quoted wrong-key
classification values are also withheld because their scaling
artifact is not checked in.  A nonzero signal
at an excluded coordinate contradicts the construction; a nonzero
end-to-end attack metric can also reveal an implementation error or
a path outside the declared isolation surface.

Prior work~\citep{mccormick2026lora} reports end-to-end leakage
measurements using a gated control.  Those external attack studies
are not reproduced here and are not used to enlarge the local FFN
theorem.  This paper's empirical evidence is limited to
\Cref{sec:empirical} and \Cref{sec:superposition-empirical}.

\paragraph{Governance relevance.}
The runtime binds an authenticated grant, model version, derived
mask, and audit record before applying the gate.  This can support
provenance questions such as which FFN coordinates were authorized
for an invocation and whether a confined slice was revoked.  It does
not by itself satisfy a regulatory framework, certify shared model
components, or prove deletion of state outside the declared
parameter and runtime surface.  Nor does deleting a recipient from a
new lockbox revoke copied key bytes unless the runtime rejects the old
grant ID/key epoch.  Under that mandatory online check,
principal-scoped derivation permits one tenant + sub-id grant to be
revoked without rotating unrelated principals.

\paragraph{Comparison with LoRA isolation.}
LoRA does not itself provide a cryptographic authorization boundary.
CDP structurally eliminates any component of a signal path that must
traverse a non-owning gated FFN coordinate; declared shared
components remain outside that statement.
The relevant distinction is architectural: LoRA adds trainable
low-rank updates but does not impose an authorization gate, whereas
CDP's local statement follows from Hadamard zeroing at the declared
FFN surface.

\paragraph{Limitations.}
For intra-model FFN partitioning, the primary limitation is the shared
backbone trust boundary (\Cref{sec:trust-boundary}).  Attention,
embeddings, and the output projection are shared.  We do not evaluate
leakage through those shared activation paths.  The proofs and
exact-delta experiments cover the declared FFN/private state boundary,
not partitioning of the complete model.  Whole-model gating is a
separate operational authorization boundary and does not strengthen
the local FFN theorem.

The guarantee also depends on gate placement and path closure.  A
post-mixing gate or an ungated bypass may preserve the local identity
$0\cdot x=0$ while invalidating the intended end-to-end isolation
surface.  Implementations must place the FFN gate on the expanded
activation before the down projection for the intermediate theorem,
or before the first projection for the narrower pre-MLP surface.
Execution-scoped state must
follow the same regime boundary wherever it is claimed as confined.
Attention-coordinate partitioning and generative tenant-state
training remain outside the evaluated scope.  Cargo's frozen-generator
smoke tests runtime release of retrieved context, not deposition into
generative model weights.

The formative five-seed classification sweep gates the final pooled
representation rather than the expanded intermediate FFN activation.
It therefore demonstrates support-constrained task feasibility at
that downstream surface.  The separate strict-cotenancy study does
place the gate at every expanded intermediate FFN and reports utility
and exact parameter/optimizer confinement through $R{=}16$.
Neither study partitions shared attention coordinates or establishes
multi-regime isolation across the internals of a complete model.

The strict experiments use one AG News distribution and label
permutations as regime-specific functional identities.  This is
sufficient to exercise routing, aligned updates, and exact inactive
state checks, but it is weak evidence for heterogeneous customer
corpora or naturally different tenant tasks.  A stronger follow-up
should use distinct corpora or task distributions with shared label
semantics, paired seeds, and a deliberately broken gate or bypass
control.

\paragraph{Broader impact.}
CDP can support consent-based policy for the declared gated FFN surface.
Coordinate exclusion is the default state there; intentional unions
require an authorized multi-regime grant.  Shared components remain
inside the stated trust boundary.  The mechanism may be useful where
operators need an auditable local state boundary, but it does not
itself establish regulatory compliance or an end-to-end privacy
guarantee.

Algebraic annihilation (\Cref{sec:destruction}) provides verifiable
FFN-slice erasure: the selected $\bW_1$, $\mathbf{b}_1$, and $\bW_2$
coordinates are zeroed, and the confinement theorem states that
regime-scoped updates did not leave that declared parameter surface.
Calling this a complete right-to-deletion mechanism additionally
requires $\mathbf{b}_2$, optimizer state, adapters, shared
components, logs, caches, and backups to be frozen, separately
scoped, or erased under the same lifecycle policy.

% ===================================================================
\section{Conclusion}
\label{sec:conclusion}

Cryptographic Dimension Partitioning defines two explicit scopes.
Within an FFN, a globally assigned binary mask provides exact
coordinate zeros at either a pre-MLP input or an expanded intermediate
activation; the strongest aligned parameter-state result uses the
intermediate placement.  Around a complete model, an IdP-bound runtime
gate authorizes or rejects the unchanged model as one atomic resource.
The core Lean~4 theorems contain no \texttt{sorry} declarations or
project-specific axioms.  Their scope is precise: exact forward and
loss-gradient zeros at a correctly placed FFN gate, and aligned
parameter-state confinement when the optimizer obeys the same support.

The strict DistilBERT study measures the resulting tradeoff.  Across
three seeds, owning accuracy falls from $91.28\%$ at $R{=}1$ to
$86.21\%$ at $R{=}16$, while frozen-shared, off-partition parameter,
off-partition optimizer-moment, and inactive-classifier deltas remain
exactly zero.  Private residual adapters and private experts recover
$90.47\%$ and $90.09\%$ mean accuracy without modifying the shared
backbone or inactive lanes.  The separate post-encoder study shows
small support-constrained classification costs, but does not validate
intermediate-FFN placement. The equal-work public-adaptation factorial separates
ordinary extra training, mask-aware adaptation, and distillation conditional on
masks; its positive result is one full seed and therefore not a
population-level effect estimate.

The authority-constrained MoE study adds a distinct result: learned semantic
routing can operate inside a fixed grant-selected expert set. At $R{=}8$,
packing separately trained routers and experts changes zero held-out logits or
predictions while unauthorized dispatch and the audited inactive training
state remain exactly zero. This is lossless module packing, not compression of
expert bytes or partitioning of shared attention.

The capability-prefix extension sharpens the control/data boundary at token
granularity. Execution authority, not prompt text, selects the private continuous
prefix and allowed experts; the learned token router operates only afterward.
The held-out spoof control and 311,477 routed-token trace support that bounded
claim, but the classifier does not establish autoregressive generation quality.

An intra-model FFN gate does not partition attention, embeddings,
normalization, output heads, caches, or undeclared bypasses.  The
whole-model gate covers authorization to invoke those components as an
atomic unit but does not partition them internally.  Cargo
authorization and keyed permutation conjugation are supplementary
runtime and reparameterization studies, not extensions of the FFN
theorem.
Generative tenant-state training, heterogeneous private-corpus
validation, and multi-regime internal isolation of a complete shared
model remain open.  The
paper, Lean proofs, source-visible research preflight, experiment code, and
artifacts are available at
\url{https://github.com/sekosai/schemen-gate/tree/main/research/cdp}.
Live identity, key-management, model-serving, network, and audit
integrations require separately pinned deployment evidence.

% ===================================================================
\bibliographystyle{plainnat}

\begin{thebibliography}{27}

\bibitem[Abadi et~al.(2016)]{abadi2016deep}
M.~Abadi, A.~Chu, I.~Goodfellow, H.~B. McMahan, I.~Mironov, K.~Talwar, and
  L.~Zhang.
\newblock Deep learning with differential privacy.
\newblock In \emph{CCS}, 2016.
\newblock \url{https://arxiv.org/abs/1607.00133}

\bibitem[Chen et~al.(2023)]{chen2023punica}
L.~Chen, Z.~Ye, Y.~Wu, D.~Zhuo, L.~Ceze, and A.~Krishnamurthy.
\newblock Punica: Multi-tenant LoRA serving.
\newblock \emph{MLSys}, 2024.
\newblock \url{https://arxiv.org/abs/2310.18547}

\bibitem[Cohen et~al.(2019)]{cohen2019certified}
J.~Cohen, E.~Rosenfeld, and Z.~Kolter.
\newblock Certified adversarial robustness via randomized smoothing.
\newblock In \emph{ICML}, 2019.
\newblock \url{https://arxiv.org/abs/1902.02918}

\bibitem[Dettmers et~al.(2022)]{dettmers2022llmint8}
T.~Dettmers, M.~Lewis, Y.~Belkada, and L.~Zettlemoyer.
\newblock LLM.int8(): 8-bit matrix multiplication for transformers at scale.
\newblock In \emph{NeurIPS}, 2022.
\newblock \url{https://arxiv.org/abs/2208.07339}

\bibitem[Dettmers et~al.(2024)]{dettmers2024qlora}
T.~Dettmers, A.~Pagnoni, A.~Holtzman, and L.~Zettlemoyer.
\newblock QLoRA: Efficient finetuning of quantized language models.
\newblock In \emph{NeurIPS}, 2024.
\newblock \url{https://arxiv.org/abs/2305.14314}

\bibitem[Elhage et~al.(2022)]{elhage2022superposition}
N.~Elhage, T.~Hume, C.~Olsson, et~al.
\newblock Toy models of superposition.
\newblock \emph{Transformer Circuits Thread}, 2022.
\newblock \url{https://transformer-circuits.pub/2022/toy_model/index.html}

\bibitem[Fedus et~al.(2022)]{fedus2022switch}
W.~Fedus, B.~Zoph, and N.~Shazeer.
\newblock Switch transformers: Scaling to trillion parameter models.
\newblock \emph{JMLR}, 2022.
\newblock \url{https://arxiv.org/abs/2101.03961}

\bibitem[Frantar et~al.(2023)]{frantar2023gptq}
E.~Frantar, S.~Ashkboos, T.~Hoefler, and D.~Alistarh.
\newblock GPTQ: Accurate post-training quantization for generative
  pre-trained transformers.
\newblock In \emph{ICLR}, 2023.
\newblock \url{https://arxiv.org/abs/2210.17323}

\bibitem[Hu et~al.(2022)]{hu2022lora}
E.~J. Hu, Y.~Shen, P.~Wallis, Z.~Allen-Zhu, Y.~Li, S.~Wang, L.~Wang, and
  W.~Chen.
\newblock LoRA: Low-rank adaptation of large language models.
\newblock In \emph{ICLR}, 2022.
\newblock \url{https://arxiv.org/abs/2106.09685}

\bibitem[Kwon et~al.(2023)]{kwon2023vllm}
W.~Kwon, Z.~Li, S.~Zhuang, Y.~Sheng, L.~Zheng, C.~H. Yu, J.~Gonzalez,
  H.~Zhang, and I.~Stoica.
\newblock Efficient memory management for large language model serving with PagedAttention.
\newblock In \emph{SOSP}, 2023.
\newblock Apache-2.0 license.
\newblock \url{https://arxiv.org/abs/2309.06180}

\bibitem[Lepikhin et~al.(2021)]{lepikhin2021gshard}
D.~Lepikhin, H.~Lee, Y.~Xu, D.~Chen, O.~Firat, Y.~Huang, M.~Krikun,
  N.~Shazeer, and Z.~Chen.
\newblock GShard: Scaling giant models with conditional computation and
  automatic sharding.
\newblock In \emph{ICLR}, 2021.
\newblock \url{https://arxiv.org/abs/2006.16668}

\bibitem[Li and Liang(2021)]{li2021prefix}
X.~L. Li and P.~Liang.
\newblock Prefix-tuning: Optimizing continuous prompts for generation.
\newblock In \emph{ACL-IJCNLP}, 2021.
\newblock \url{https://aclanthology.org/2021.acl-long.353/}

\bibitem[Ma et~al.(2024)]{ma2024bitnet158}
S.~Ma, H.~Wang, L.~Ma, et~al.
\newblock The era of 1-bit LLMs: All large language models are in 1.58 bits.
\newblock \emph{arXiv:2402.17764}, 2024.
\newblock \url{https://arxiv.org/abs/2402.17764}

\bibitem[Mallya and Lazebnik(2018)]{mallya2018packnet}
A.~Mallya and S.~Lazebnik.
\newblock PackNet: Adding multiple tasks to a single network by iterative pruning.
\newblock In \emph{CVPR}, 2018.
\newblock \url{https://arxiv.org/abs/1711.05769}

\bibitem[Mallya et~al.(2018)]{mallya2018piggyback}
A.~Mallya, D.~Davis, and S.~Lazebnik.
\newblock Piggyback: Adapting a single network to multiple tasks by learning to mask weights.
\newblock In \emph{ECCV}, 2018.
\newblock \url{https://arxiv.org/abs/1801.06519}

\bibitem[McCormick(2026)]{mccormick2026lora}
R.~McCormick.
\newblock Multi-tenant LoRA: Attention is all you bleed.
\newblock \emph{In preparation}, 2026.

\bibitem[Predibase(2024)]{predibase2024lorax}
Predibase.
\newblock LoRAX: Multi-LoRA inference server.
\newblock Apache-2.0 license.
\newblock \url{https://github.com/predibase/lorax}, 2024.

\bibitem[Sanh et~al.(2019)]{sanh2019distilbert}
V.~Sanh, L.~Debut, J.~Chaumond, and T.~Wolf.
\newblock DistilBERT, a distilled version of BERT: smaller, faster, cheaper and lighter.
\newblock \emph{arXiv:1910.01108}, 2019.
\newblock Apache-2.0 license.
\newblock \url{https://arxiv.org/abs/1910.01108}

\bibitem[Seshia et~al.(2018)]{seshia2018formal}
S.~A. Seshia, A.~Desai, T.~Dreossi, D.~J. Fremont, S.~Ghosh,
  E.~Kim, S.~Shivakumar, M.~Vazquez-Chanlatte, and X.~Yue.
\newblock Formal specification for deep neural networks.
\newblock In \emph{ATVA}, 2018.
\newblock \url{https://arxiv.org/abs/1710.01688}

\bibitem[Sheng et~al.(2024)]{sheng2024slora}
Y.~Sheng, S.~Cao, D.~Li, et~al.
\newblock S-LoRA: Serving thousands of concurrent LoRA adapters.
\newblock In \emph{MLSys}, 2024.
\newblock \url{https://arxiv.org/abs/2311.03285}

\bibitem[SPIFFE Project(2026)]{spiffe2026}
SPIFFE Project.
\newblock SPIFFE Workload API specification.
\newblock Cloud Native Computing Foundation, 2026.
\newblock \url{https://spiffe.io/docs/latest/spiffe-specs/spiffe_workload_api/}

\bibitem[Tram\`er and Boneh(2018)]{tramer2018slalom}
F.~Tram\`er and D.~Boneh.
\newblock Slalom: Fast, verifiable and private execution of neural networks in trusted hardware.
\newblock In \emph{ICLR}, 2019.
\newblock \url{https://arxiv.org/abs/1806.03287}

\bibitem[Wang and Li(2023)]{wang2023dpfederated}
Z.~Wang and J.~Li.
\newblock Differentially private federated learning with LoRA.
\newblock \emph{arXiv:2306.05908}, 2023.
\newblock \url{https://arxiv.org/abs/2306.05908}

\bibitem[Wortsman et~al.(2020)]{wortsman2020supermasks}
M.~Wortsman, V.~Ramanujan, R.~Liu, A.~Kembhavi, M.~Rastegari,
J.~Yosinski, and A.~Farhadi.
\newblock Supermasks in superposition.
\newblock In \emph{NeurIPS}, 2020.
\newblock \url{https://arxiv.org/abs/2006.14769}

\bibitem[Zhang et~al.(2015)]{zhang2015agnews}
X.~Zhang, J.~Zhao, and Y.~LeCun.
\newblock Character-level convolutional networks for text classification.
\newblock In \emph{NeurIPS}, 2015.
\newblock Non-commercial research license.
\newblock \url{https://arxiv.org/abs/1509.01626}

\end{thebibliography}

\end{document}
